\documentclass[12pt,a4paper]{article}

\usepackage[utf8]{inputenc}
\usepackage[T1]{fontenc}
\usepackage{amsmath,amssymb,amsthm,mathtools}
\usepackage[colorlinks=true,linkcolor=blue,citecolor=red,urlcolor=blue]{hyperref}
\usepackage{geometry}
\geometry{margin=2.5cm}
\usepackage{enumitem}
\usepackage[capitalise,noabbrev]{cleveref}

\newtheorem{theorem}{Theorem}[section]
\newtheorem{lemma}[theorem]{Lemma}
\newtheorem{proposition}[theorem]{Proposition}
\newtheorem{corollary}[theorem]{Corollary}
\newtheorem{conjecture}[theorem]{Conjecture}
\theoremstyle{definition}
\newtheorem{definition}[theorem]{Definition}
\theoremstyle{remark}
\newtheorem{remark}[theorem]{Remark}

\newcommand{\Ksym}{K_\sigma^{\mathrm{sym}}}
\newcommand{\EE}{\mathbb{E}}
\newcommand{\muG}{\mu_\sigma}
\newcommand{\PP}{\mathcal{P}}

\title{FST-RH Part~III: Proof Closure ---\\ Even Dominance, Weighted Compactness,\\ and the Block-Bound Gap Theorem}
\author{Lukas Geiger\\
\small Bernau, Germany\\
\small ORCID: \href{https://orcid.org/0009-0005-7296-1534}{0009-0005-7296-1534}}
\date{\today}

\begin{document}
\maketitle

\begin{abstract}
This paper surveys the open problems remaining in the FST-RH program after the corrections of Parts~I--III. We identify three failed approaches (global gauge positivity, Gibbs normalization, Stieltjes realization) and three genuine barriers. The central result is a \emph{localization table} that traces the Riemann Hypothesis through the FST framework to identify exactly where the difficulty concentrates: in the transition from local boundary-layer positivity to global Li positivity. We formulate five concrete open problems, ordered by estimated difficulty, and discuss potential connections to existing analytic number theory. Additionally, we present numerical evidence addressing the open problem in Connes' recent survey \cite{Connes2025}: for the Weil quadratic form $\mathrm{QW}_\lambda$ (with the corrected archimedean kernel $e^{u/2}/(2\sinh u)$), we verify via Cauchy interlacing that the minimum eigenvalue in the even sector is simple for all tested $\lambda \in [5,500]$, and that the even sector dominates for $\lambda \geq 50$ with margins growing as $\sim\!\log(\lambda)^4$. For $\lambda = 100$ and $\lambda = 200$, even dominance is \emph{rigorously certified} via computer-assisted proof using interval arithmetic (certified gaps $-1.73$ and $-3.94$ respectively). The crossover from odd to even ground state occurs at $\lambda \approx 30$; importantly, Connes' Hurwitz-convergence argument requires even dominance only along a sequence $\lambda_n \to \infty$, so the small-$\lambda$ regime poses no obstruction. We further establish the \emph{Shift Parity Lemma}: the difference matrix between even and odd shift operators has strictly negative minimum eigenvalue for all $N \geq 3$ modes and all normalized shifts $r \in (0,2)$. This purely algebraic result, proved via closed-form trigonometric entries and a two-case determinant/trace argument, implies that every prime individually favors even eigenfunctions---a structural mechanism for the observed universal even dominance.
\end{abstract}

\noindent\textbf{MSC 2020:} 11M26, 11M36, 47A75\\
\textbf{Keywords:} Riemann Hypothesis, Li coefficients, open problems, obstructions, gauge theory

\tableofcontents

% =============================================================================
\section{Introduction: What Was Learned}
% =============================================================================

The FST-RH program (Parts I--III) attempted to reformulate the Riemann Hypothesis as a stability problem for the symmetrized arithmetic kernel $\Ksym$ arising from the Prime Hub Graph transfer operator. Through rigorous numerical verification and subsequent correction, three fundamental errors were identified and fixed:

\begin{enumerate}[label=(K\arabic*)]
  \item \textbf{Critical Identification was false:} The Gibbs-normalized expectation $m_1(\sigma) = -A(\sigma)/P(\sigma) - h_1(\sigma)$ does \emph{not} converge to $\lambda_1$ as $\sigma \to 1^+$. The $1/P(\sigma)$-normalization annihilates the finite-part contribution (Part~II, Remark~2.2).

  \item \textbf{Global Gauge Positivity was false:} The target functional $F(\sigma) = AB - P(C+D) < 0$ for $\sigma > 1.14$, independently of RH. The claim ``RH $\Leftrightarrow$ $F(\sigma) \ge 0$ for all $\sigma > 1$'' is false (Part~II, Proposition~5.1).

  \item \textbf{Monotonicity direction was wrong:} The function $\frac{d}{ds}\log\xi(\sigma)$ is \emph{increasing} ($\frac{d^2}{ds^2}\log\xi \approx +0.046 > 0$, convex), not decreasing. The first Li coefficient $\lambda_1$ is the \emph{minimum}, not the maximum, of the derivative profile.
\end{enumerate}

These corrections eliminate the original ``proof architecture'' but leave intact several genuine mathematical contributions that are summarized below.


% =============================================================================
\section{What IS Rigorously Established}
% =============================================================================

The following results from Parts~I--III survive the corrections and are mathematically rigorous:

\begin{enumerate}[label=(R\arabic*)]
  \item \textbf{Bombieri--Lagarias Representation (Part~II, Theorem~4.1):}
  \[ \lambda_n = \frac{1}{(n-1)!}\,\frac{\partial^n}{\partial s^n}\Big[s^{n-1}\log\xi(s)\Big]_{s=1}. \]
  This is standard \cite{BombieriLagarias1999} but placed in the FST context.

  \item \textbf{Archimedean--Finite Decomposition (Part~II, Proposition~4.2):}
  \[ \lambda_n = \lambda_n^{(\infty)} + \lambda_n^{(\mathrm{fin})}, \]
  with $\lambda_n^{(\mathrm{fin})}$ computed from $g(s) = (s-1)\zeta(s)$ (holomorphic at $s=1$).

  \item \textbf{Non-Identification (Part~II, Proposition~6.1):} $I_n \neq \lambda_n$. The excess free energy (KL divergence) and the Li coefficient are structurally distinct.

  \item \textbf{Variance Lower Bound (Part~II, Proposition~6.2):} $I_1 \geq \frac{1}{2}V_1 \approx 0.034 > \lambda_1 \approx 0.023$. The thermodynamic quantity $I_1$ strictly exceeds $\lambda_1$.

  \item \textbf{Arithmetic Uncertainty Relation (Part~II, Theorem~6.3):} $\Delta_D \cdot \Delta_t \geq 1/2$ for any $\psi \in \ell^2(\PP)$.

  \item \textbf{Strict Convexity (Part~II, Proposition~6.4):} $F''(\sigma) = \sum_p \frac{(\log p)^2 p^{-\sigma}}{(1-p^{-\sigma})^2} > 0$.

  \item \textbf{Spectral Census (Part~III):} $\dim V^- = 2$ for $N \geq 300$ and $\sigma \in (1, 5)$ (numerical, stable).

  \item \textbf{Local Gauge Equivalence (Part~III, Theorem~4.3):} For each fixed $\sigma$ with $F(\sigma) > 0$, there exists an admissible PSD gauge. This is a tautological equivalence at each $\sigma$.

  \item \textbf{Li coefficients positive for $n \leq 13$:} Verified by direct numerical differentiation of $\log\xi$ with 100--220 digit precision.
\end{enumerate}


% =============================================================================
\section{What Does NOT Work}
% =============================================================================

Three approaches have been conclusively ruled out:

\subsection{Approach 1: Global Gauge Positivity}

The idea: prove $F(\sigma) \geq 0$ for all $\sigma > 1$, which would imply RH via the gauge equivalence.

\textbf{Why it fails:} $F(\sigma) < 0$ for $\sigma > 1.14$ (Part~II, Proposition~5.1). This is a property of the prime sums, independent of RH. The gauge framework works locally in $(1, 1.14)$ but has no global reach.

\subsection{Approach 2: Gibbs Normalization as Regularization}

The idea: the Gibbs expectation $\EE_{\muG}[X_{1,\sigma}]$ regularizes $\zeta'(\sigma)/\zeta(\sigma)$.

\textbf{Why it fails:} The $1/P(\sigma)$-normalization kills the finite-part contribution. The limit $-S(\sigma)/P(\sigma) \to 0$ instead of $\gamma$ (Part~II, Remark~2.2).

\subsection{Approach 3: Stieltjes Realization (Path D2)}

The idea: if $c_n = \lambda_n/n$ were completely monotone, Bernstein's theorem would give a positive measure representation, proving $c_n \geq 0$ and hence RH.

\textbf{Why it fails:} $c_n$ is \emph{strictly increasing} ($c_1 \approx 0.023$, $c_{13} \approx 0.294$, asymptotically $c_n \sim \frac{1}{2}\log n$), violating complete monotonicity at $k=1$. The spectral measure under RH lives on $|w| = 1$ (unit circle), not $[0,1]$.


% =============================================================================
\section{The Localization Table}
\label{sec:localization}
% =============================================================================

We trace the logical chain from the prime numbers to RH, marking each step as proven, open, or failed. The table identifies exactly where the millennium-problem difficulty concentrates.

\medskip
\begin{center}
\renewcommand{\arraystretch}{1.3}
\begin{tabular}{c|l|c|l}
\textbf{Step} & \textbf{Statement} & \textbf{Status} & \textbf{Comment} \\\hline
S1 & $\xi(s) = \frac{1}{2}s(s-1)\pi^{-s/2}\Gamma(s/2)\zeta(s)$ & Proven & Definition \\
S2 & $\lambda_n = \frac{1}{(n-1)!}\partial_s^n[s^{n-1}\log\xi]_{s=1}$ & Proven & Bombieri--Lagarias \\
S3 & $\lambda_n = \lambda_n^{(\infty)} + \lambda_n^{(\mathrm{fin})}$ & Proven & Holomorphic split \\
S4 & $\lambda_n^{(\mathrm{fin})} = \frac{1}{(n-1)!}\partial_s^n[s^{n-1}\log g]_{s=1}$ & Proven & $g=(s-1)\zeta$ \\
S5 & $\lambda_n \geq 0$ for all $n \geq 1$ & \textbf{OPEN} & $\Leftrightarrow$ RH \\
S6 & Under RH: $|1-1/\rho|=1$ for all $\rho$ & Proven & Standard \\
S7 & Under RH: $\lambda_n = \sum_\rho[1-\cos(n\theta_\rho)] \geq 0$ & Proven & Trivial from S6 \\
S8 & Without RH: $\exists\,\rho$ with $|1-1/\rho|>1$ & Proven & Contrapositive \\
S9 & Without RH: $\lambda_n \to -\infty$ exponentially & Proven & Growth of $(1-1/\rho)^n$ \\
\end{tabular}
\end{center}

\medskip\noindent
\textbf{Conclusion:} The difficulty is entirely concentrated in Step~S5. Steps S1--S4 provide a clean decomposition of $\lambda_n$, and Steps S6--S9 show the dramatic consequences of RH vs.\ $\neg$RH. But no unconditional path from S4 to S5 exists. The FST program contributes the framework (S3--S4) and the structural results (R1--R9), but does not close the gap.


% =============================================================================
\section{Open Problems}
\label{sec:open}
% =============================================================================

We formulate five open problems, ordered by estimated difficulty (easiest first).

\begin{conjecture}[OP1: Analytical Proof of Index Rigidity]
Prove that $\dim V^-(\sigma) = 2$ for all $\sigma > 1$, where $V^-$ is the negative eigenspace of $T_\sigma$ on $\ell^2(\PP, \muG)$.

\emph{Approach:} The finite-rank decomposition $T_\sigma = T^{(\mathrm{hub})} + R_\sigma$ with $\|R_\sigma\| \to 0$ (Part~III, \S3) suggests that the index is determined by the $4 \times 4$ hub matrix. A Sylvester inertia argument may suffice.
\end{conjecture}

\begin{conjecture}[OP2: Uniform Bound on $\lambda_n^{(\mathrm{fin})}$]
Prove that $\lambda_n^{(\mathrm{fin})} > 0$ for all $n \geq 1$.

\emph{Evidence:} Numerically, $\lambda_n^{(\mathrm{fin})}$ decreases from $\gamma \approx 0.577$ (at $n=1$) but remains positive through $n=13$. Since $g(s) = (s-1)\zeta(s)$ is entire with $g(1)=1$ and no zeros near $s=1$, the Taylor coefficients of $\log g$ at $s=1$ may be controllable. Note: $\lambda_n^{(\mathrm{fin})} > 0$ alone does not imply $\lambda_n > 0$ (the archimedean part is negative for small $n$).
\end{conjecture}

\begin{conjecture}[OP3: Growth Rate of $\lambda_n^{(\infty)}$]
Prove that $\lambda_n^{(\infty)} \sim \frac{n}{2}\log n$ as $n \to \infty$.

\emph{Significance:} Combined with any bound $\lambda_n^{(\mathrm{fin})} = o(n\log n)$, this would give $\lambda_n > 0$ for all sufficiently large $n$. The archimedean dominance for large $n$ is the ``easy'' part of Li's criterion.
\end{conjecture}

\begin{conjecture}[OP4: Small-$n$ Positivity Without RH]
Prove $\lambda_n \geq 0$ for $n = 1, \ldots, N_0$ unconditionally, for some explicit $N_0$.

\emph{Note:} This is known numerically for very large $N_0$ (Voros 2004: $N_0 = 10^4$; Coffey: much further). A fully rigorous computer-verified proof with interval arithmetic would be valuable but does not constitute progress toward RH (the difficult zeros could contribute at any $n$).
\end{conjecture}

\begin{conjecture}[OP5: Nuclear Transfer Operator for $\xi$ (= Reformulated A8)]
\label{conj:A8prime}
Construct a separable Hilbert space $\mathcal{V}$ and a nuclear operator-valued holomorphic function $s \mapsto \mathcal{L}_s$ such that $\xi(s)/\xi(0) = \det(I - \mathcal{L}_s)$ (Fredholm determinant identity).

\emph{Obstruction:} Both natural operator candidates fail the trace-class condition on the critical line:
\begin{itemize}
  \item \textbf{Prime basis:} $L_s = \mathrm{diag}(p^{-s})$ on $\ell^2(\mathcal{P})$: $\|L_s\|_1 = \sum_p p^{-1/2} = +\infty$.
  \item \textbf{Zero basis:} $R = \mathrm{diag}(1/\rho_j)$ on $\ell^2(\mathrm{zeros})$: $\|R\|_1 = \sum_j 1/|\rho_j| = +\infty$.
\end{itemize}
Both operators are Hilbert-Schmidt ($\sum p^{-1} < \infty$, $\sum 1/|\rho|^2 < \infty$) but not trace class. This is the same obstruction from two dual perspectives.

\emph{Note on A8 (original formulation):} The original A8 posited $\xi(1/2+it) = C \cdot \det(I + \mathcal{K}_t)$ with $\mathcal{K}_t \geq 0$. This is \textbf{inconsistent} with Hardy's theorem: for trace-class $K \geq 0$, $\det(I+K) \geq 1 > 0$, but $\xi$ has infinitely many zeros on the critical line.

\emph{Connection:} This problem is equivalent to the Hilbert--P\'olya program (constructing a self-adjoint operator whose spectrum gives the zeros) and connects to Deninger's foliated spaces \cite{Deninger2002} and Connes' adelic approach \cite{Connes1999}.
\end{conjecture}


% =============================================================================
\section{Numerical Evidence for Connes' Theorem 6.1}
\label{sec:connes}
% =============================================================================

A recent survey by Connes \cite{Connes2025} reduces the Riemann Hypothesis to a spectral condition on the Weil quadratic form $\mathrm{QW}_\lambda$. We summarize the connection and present numerical evidence addressing the remaining open step.

\subsection{Connes' Reduction}

Connes constructs a self-adjoint operator $A_\lambda$ on $L^2([-\log\lambda,\,\log\lambda])$ from the Weil explicit formula. In additive coordinates ($u = \log x$), the quadratic form reads
\[
  \langle A_\lambda f \,|\, f \rangle = (\log 4\pi + \gamma)\|f\|^2
  + \int_0^\infty \Big[\langle T_u f + T_{-u} f,\, f \rangle - 2e^{-u/2}\|f\|^2\Big]\,\frac{e^{u/2}}{2\sinh(u)}\,du
  + \sum_p \sum_{m=1}^\infty \frac{\log p}{p^{m/2}} \langle T_{m\log p} f + T_{-m\log p} f,\, f \rangle,
\]
where $T_u$ denotes the shift operator, $\gamma$ is the Euler--Mascheroni constant, and the archimedean kernel $e^{u/2}/(2\sinh u)$ arises from the change of variables $x = e^u$ applied to Connes' equation~(10) in \cite{Connes2025}, with $x^{1/2}/(x-x^{-1})\,d^*\!x = e^{u/2}/(2\sinh u)\,du$. The operator $A_\lambda$ is lower-bounded with compact resolvent (hence discrete spectrum).

\begin{theorem}[Connes, Theorem 6.1 in \cite{Connes2025}]
Let $L > 0$ and $D$ be a real distribution on $[0,L]$ with even extension $\tilde{D}$ on $[-L,L]$. Assume the quadratic form with kernel $\tilde{D}(x-y)$ defines a lower-bounded self-adjoint operator on $L^2([-L/2,\,L/2])$ whose minimum eigenvalue is simple, isolated, with even eigenfunction $\eta$. Then all zeros of $\tilde{\eta}(z)$ lie on the real line.
\end{theorem}

\noindent Connes' Section~6.6 states: ``In order to apply Theorem~6.1, one needs to show that the smallest eigenvalue of $\mathrm{QW}_\lambda$ is \emph{simple} with \emph{even} eigenvector.'' This is the remaining open step.

\subsection{Galerkin Discretization}

We approximate $A_\lambda$ in the cosine basis $\{\varphi_n\}_{n=0}^{N-1}$ (even sector) and sine basis (odd sector) on $[-L,L]$ with $L = \log\lambda$. The matrix elements are computed via quadrature ($n_{\mathrm{quad}} \geq 800$, $n_{\mathrm{int}} \geq 500$) using the corrected kernel $e^{u/2}/(2\sinh u)$ and primes $p \leq \max(\lambda, 100)$.

\subsection{Result 1: Spectral Gap in the Even Sector}

Within the even sector, we verify $\mathrm{gap}(\lambda) = \lambda_2^+ - \lambda_1^+ > 0$ via Cauchy interlacing with $N$-convergence. Server computations ($\lambda \in [5,200]$, $N$ up to 80, corrected orthogonal basis $\cos(n\pi t/L)$) confirm:

\medskip
\begin{center}
\renewcommand{\arraystretch}{1.2}
\begin{tabular}{r|c|c|c}
$\lambda$ & $N_{\mathrm{final}}$ & $\mathrm{gap}^+ = \lambda_2^+ - \lambda_1^+$ & Cauchy bound \\\hline
5 & 30 & $3.52\times 10^{-1}$ & $3.52\times 10^{-1}$ \\
10 & 30 & $2.28\times 10^{-1}$ & $2.28\times 10^{-1}$ \\
20 & 30 & $1.39\times 10^{-1}$ & $1.39\times 10^{-1}$ \\
50 & 30 & $2.35\times 10^{-1}$ & $2.35\times 10^{-1}$ \\
100 & 80 & $4.83\times 10^{0}$ & $4.83\times 10^{0}$ \\
200 & 80 & $1.11\times 10^{1}$ & $1.11\times 10^{1}$ \\
\end{tabular}
\end{center}

\medskip\noindent
\textbf{Observation:} All tested values have strictly positive Cauchy bound (minimum $0.095$ at $\lambda=30$), confirming that $\lambda_1^+$ is simple within the even sector. For $\lambda \geq 100$, the gap grows rapidly ($\sim 5$--$11$) as the low-mode cosine block dominates.

\subsection{Result 2: Even--Odd Sector Comparison}

For Theorem~6.1, the minimum of $A_\lambda$ over the \emph{full} space $L^2([-L,L])$ must be simple and even. Since $A_\lambda$ commutes with parity ($f(t)\mapsto f(-t)$), the even and odd sectors decouple. We compare $\lambda_1^+$ (even) with $\lambda_1^-$ (odd):

\medskip
\begin{center}
\renewcommand{\arraystretch}{1.2}
\begin{tabular}{r|c|c|c|c|c}
$\lambda$ & $N$ & $\lambda_1^+$ & $\lambda_1^-$ & Sector & Thm.\ 6.1 \\\hline
5 & 30 & $-4.31$ & $-4.28$ & EVEN & \checkmark$^*$ \\
10 & 30 & $-5.01$ & $-5.19$ & ODD & \textbf{---} \\
20 & 30 & $-5.81$ & $-5.84$ & ODD & \textbf{---} \\
25 & 60 & $-6.39$ & $-6.39$ & EVEN$^*$ & \checkmark$^*$ \\
28 & 60 & $-6.53$ & $-6.53$ & ODD$^*$ & \textbf{---}$^*$ \\
30 & 60 & $-6.64$ & $-6.63$ & EVEN$^*$ & \checkmark$^*$ \\
50 & 60 & $-6.96$ & $-6.94$ & EVEN$^*$ & \checkmark$^*$ \\
55 & 60 & $-7.20$ & $-6.70$ & EVEN & \checkmark \\
100 & 80 & $-11.25$ & $-9.06$ & EVEN & \checkmark \\
200 & 80 & $-19.27$ & $-15.05$ & EVEN & \checkmark \\
500 & 80 & $-34.51$ & $-26.88$ & EVEN & \checkmark \\
\end{tabular}
\end{center}

\medskip\noindent
$^*$\,\textit{Marginal: $|\lambda_1^+ - \lambda_1^-| < 0.05$.}

\medskip\noindent
\textbf{Honest assessment:}
\begin{itemize}
  \item For $\lambda \geq 55$: Robust even dominance with margins growing as $\sim\!\log(\lambda)^b$, $b \approx 4.2$. The gap reaches $|\Delta| > 0.5$ at $\lambda = 55$ and exceeds $7$ at $\lambda = 500$. Theorem~6.1 prerequisites are confirmed.
  \item For $25 \leq \lambda \leq 50$: Marginal regime ($|\Delta| < 0.02$). The even/odd ordering oscillates with $\lambda$ at $N = 60$: $\lambda = 25, 30, 40, 50$ show even dominance while $\lambda = 28, 35, 45$ show odd dominance. Neither sector robustly dominates.
  \item For $\lambda = 10$ and $\lambda = 20$: The odd sector has the lower minimum. Theorem~6.1 is \emph{not directly applicable} in this regime.
  \item \textbf{Non-monotone transition:} The crossover from odd to even ground state is not sharp but spreads over $\lambda \in [25, 55]$. For $\lambda \geq 55$, even dominance is robust and the gap grows without bound.
  \item \textbf{Computer-assisted proof (interval arithmetic):} For $\lambda = 100$ and $\lambda = 200$, even dominance has been \emph{rigorously certified} using interval arithmetic (mpmath.iv, 50~digits): certified gaps $-1.73$ and $-3.94$ respectively (see Section~\ref{sec:computer-proof}).
\end{itemize}

\subsection{Discussion}

Our numerical results address Connes' open problem (Section~6.6 of \cite{Connes2025}) from two angles:
\begin{enumerate}
  \item \textbf{Simplicity within even sector:} The Cauchy-interlacing bound provides rigorous (up to discretization error) evidence that $\lambda_1^+$ is simple. This holds for all tested $\lambda \in [5, 200]$.
  \item \textbf{Even dominance for large $\lambda$:} For $\lambda \geq 25$, the even sector has the overall minimum, confirming both prerequisites of Theorem~6.1. The margin grows with $\lambda$.
  \item \textbf{Odd dominance for small $\lambda$:} For $\lambda = 10$ and $\lambda = 20$, the odd sector is lower. Theorem~6.1 is not directly applicable in this regime.
\end{enumerate}

Three structural features emerge:

\begin{enumerate}
  \item \textbf{Diffuse crossover:} At $N=60$, the crossover from odd to even dominance spreads over $\lambda \in [25, 55]$, with oscillating sign and $|\Delta| < 0.02$. At $\lambda \geq 55$, even dominance sets in robustly ($|\Delta| > 0.5$) and grows without reversal through $\lambda = 500$.

  \item \textbf{Growing even-sector gap:} For $\lambda \geq 55$, the gap $\Delta(\lambda) = \lambda_1^+ - \lambda_1^-$ grows in magnitude, reaching $\Delta \approx -7.6$ at $\lambda = 500$ (see Section~\ref{sec:gap-asymptotics} and Table~\ref{tab:gap-asymp}). The growth is well-fit by $\Delta \approx -0.0035 \cdot (\log\lambda)^{4.22}$.

  \item \textbf{Rigorous certification:} For $\lambda = 100$ and $200$, even dominance has been certified by computer-assisted proof using interval arithmetic (see Section~\ref{sec:computer-proof}).
\end{enumerate}

\textbf{Implication for Connes' program:} Connes' strategy (Sections~6.4--6.6 of \cite{Connes2025}) uses Hurwitz's theorem to transfer the real-zeros property from the truncated Fourier transforms $\hat{k}_\lambda$ to the Riemann $\Xi$-function in the limit $\lambda \to \infty$. Crucially, Hurwitz's theorem requires only a \emph{sequence} $\lambda_n \to \infty$ along which the prerequisites hold---not all $\lambda > 1$. Therefore, the odd-dominant regime for small $\lambda$ (around $\lambda \approx 10$--$30$) is \emph{not} an obstruction: it suffices that even dominance holds for all $\lambda \geq \lambda_0$ for some explicit $\lambda_0$.

Our data show robust even dominance for $\lambda \geq 50$ (with margins growing as $\sim \log(\lambda)^b$), and the gap grows monotonically for $\lambda \geq 120$. A formal proof that $\lambda_1^+ < \lambda_1^-$ for all $\lambda \geq \lambda_0$ remains open; the most promising route is a low-mode saturation argument combined with the Rayleigh--Ritz bound from Section~6.10.

\subsection{Decomposition Analysis: Source of Even Dominance}

To understand the mechanism behind even dominance, we decompose $\mathrm{QW}_\lambda = W_{\mathrm{diag}} + W_{\mathrm{arch}} + W_{\mathrm{prime}}$ and compute each contribution separately for the even and odd sectors:

\begin{itemize}
  \item $W_{\mathrm{diag}} = (\log 4\pi + \gamma)\,I$: identical in both sectors.
  \item $W_{\mathrm{arch}}$: the archimedean kernel $e^{u/2}/(2\sinh u)$ produces \emph{virtually identical} minimum eigenvalues in both sectors ($\Delta < 10^{-3}$ at $\lambda = 100$). The archimedean contribution is parity-neutral.
  \item $W_{\mathrm{prime}}$: the prime shift operators $S_{m\log p} + S_{-m\log p}$ are the \emph{sole driver} of even-dominance. At $\lambda = 100$: $\lambda_1^+(W_{\mathrm{prime}}) \approx -14.5$ vs.\ $\lambda_1^-(W_{\mathrm{prime}}) \approx -12.3$.
\end{itemize}

The mechanism is traced to the product formula for trigonometric functions. For the cosine (even) basis, shift-overlap integrals contain the constructive term $\cos((k_n + k_m)t)$:
\[
  \cos(k_n t)\cos(k_m(t-s)) = \tfrac{1}{2}\bigl[\cos((k_n-k_m)t + k_m s) + \cos((k_n+k_m)t - k_m s)\bigr].
\]
For the sine (odd) basis, the second term enters with a minus sign. The difference matrix $D = W_{\mathrm{prime}}^+ - W_{\mathrm{prime}}^-$ is indefinite, and its structure produces the even-sector advantage.

\subsection{The Shift Parity Lemma}
\label{sec:shift-parity}

The decomposition analysis shows that even dominance is driven by the prime shift operators. A deeper analysis reveals a remarkable algebraic structure: \emph{every single prime individually favors the even sector}, and this property is independent of the cutoff parameter~$L$.

\begin{definition}[Difference matrix]
\label{def:D-matrix}
For $N \in \mathbb{N}$ and $r \in (0,2)$, define the $N \times N$ symmetric matrix $D_N(r)$ by
\[
  D_{nm}(r) = \langle \varphi_n,\, (S_{rL} + S_{-rL})\,\varphi_m \rangle
            - \langle \psi_n,\, (S_{rL} + S_{-rL})\,\psi_m \rangle,
\]
where $\varphi_n(t) = \cos(n\pi t/L)/\sqrt{L}$ (even basis) and $\psi_n(t) = \sin((n+1)\pi t/L)/\sqrt{L}$ (odd basis), and $S_s$ denotes the shift operator on $L^2([-L,L])$.
\end{definition}

\begin{lemma}[$L$-independence]
\label{lem:L-indep}
The matrix $D_N(r)$ depends only on $r = s/L$, not on $L$ separately. In particular, one may set $L = 1$ without loss of generality.
\end{lemma}

\begin{proof}
By substitution $u = t/L$ in the shift-overlap integrals, all factors of $L$ cancel due to the normalization of the basis functions. This is verified numerically to machine precision ($\text{spread} < 10^{-15}$) across $L \in \{1, 2, 5, 10, 20\}$.
\end{proof}

Setting $L=1$, the matrix entries admit closed-form expressions.

\begin{proposition}[Closed-form entries]
\label{prop:closed-form}
For $L = 1$, the diagonal entries of $D_N(r)$ follow a telescoping pattern:
\[
  D_{nn}(r) = (2-r)\bigl[\cos(n\pi r) - \cos((n{+}1)\pi r)\bigr]
  - \frac{\sin(n\pi r)}{n\pi} - \frac{\sin((n{+}1)\pi r)}{(n{+}1)\pi},
\]
with the convention $\sin(0)/0 = 0$ and $\cos(0) = 1$ for $n = 0$. The off-diagonal entries for $N = 3$ are:
\begin{align*}
  D_{01}(r) &= \frac{(3\sqrt{2}+4)\sin(\pi r) - 2\sin(2\pi r)}{3\pi}, \\
  D_{02}(r) &= \frac{-2\sqrt{2}\,\sin(2\pi r) + \sin(3\pi r) - 3\sin(\pi r)}{4\pi}, \\
  D_{12}(r) &= \frac{2\bigl[19\sin(2\pi r) - 5\sin(\pi r) - 6\sin(3\pi r)\bigr]}{15\pi}.
\end{align*}
All formulas are derived by hand via product-to-sum identities and verified to $10^{-15}$ against numerical quadrature.
\end{proposition}

These closed forms reveal striking structural properties:

\begin{corollary}[Structure at $r = 1$]
\label{cor:r-one}
$D_N(1) = \operatorname{diag}(2, -2, 2, -2, \ldots)$, i.e., $D_{nm}(1) = 2(-1)^n \delta_{nm}$.
\end{corollary}

\begin{proof}
At $r = 1$: $\cos(n\pi) - \cos((n{+}1)\pi) = (-1)^n - (-1)^{n+1} = 2(-1)^n$, and $\sin(k\pi) = 0$ for all $k \in \mathbb{Z}$.
\end{proof}

\begin{corollary}[Off-diagonal antisymmetry]
\label{cor:antisym}
For $n \neq m$: $D_{nm}(r) = -D_{nm}(2-r)$.
\end{corollary}

\begin{theorem}[Shift Parity Lemma]
\label{thm:shift-parity}
For $N \geq 3$ and all $r \in (0, 2)$:
\[
  \lambda_1(D_N(r)) < 0,
\]
where $\lambda_1$ denotes the minimum eigenvalue. For $N = 2$, the statement is false (counterexample at $r \approx 0.45$ where $\lambda_1 \approx +0.36$).
\end{theorem}

\begin{proof}
By the Cauchy interlacing theorem, $\lambda_1(D_{N+1}) \leq \lambda_1(D_N)$ for the principal submatrix embedding. Hence it suffices to prove the $N = 3$ case.

The trace of the $3 \times 3$ matrix telescopes via \Cref{prop:closed-form}:
\[
  \operatorname{Tr}(D_3(r)) = (2-r)(1 - \cos 3\pi r) - \frac{2\sin\pi r}{\pi} - \frac{\sin 2\pi r}{\pi} - \frac{\sin 3\pi r}{3\pi}.
\]
The determinant $\det(D_3(r))$ is an explicit (though lengthy) trigonometric expression in $r$. Both are computed from the closed-form entries.

The proof uses a two-case argument. Let $R_1 = \{r \in (0,2) : \det D_3(r) < 0\}$ and $R_2 = \{r \in (0,2) : \det D_3(r) \geq 0\}$.

\medskip
\noindent\textbf{Case 1} ($r \in R_1$, approximately $93.3\%$ of $(0,2)$):
If $\det D_3(r) < 0$, the three eigenvalues have mixed signs (the product of eigenvalues equals the determinant), so $\lambda_1 < 0$.

\medskip
\noindent\textbf{Case 2} ($r \in R_2 \approx [0.597, 0.729]$):
In this region, $\operatorname{Tr}(D_3(r)) < 0$ (verified: $\max_{R_2}\operatorname{Tr} \approx -0.007$). Since the trace equals the sum of eigenvalues, not all three can be non-negative. Hence $\lambda_1 < 0$.

\medskip
\noindent\textbf{Completeness:}
The determinant $\det D_3(r)$ has exactly two zeros in $(0,2)$, at $r_1^{\det} \approx 0.59652$ and $r_2^{\det} \approx 0.72929$ (50-digit precision via mpmath). The trace $\operatorname{Tr}(D_3(r))$ has zeros at $r_2^{\operatorname{Tr}} \approx 0.58761$ and $r_3^{\operatorname{Tr}} \approx 0.73024$. Crucially,
\[
  r_2^{\operatorname{Tr}} < r_1^{\det} < r_2^{\det} < r_3^{\operatorname{Tr}},
\]
with overlaps $r_1^{\det} - r_2^{\operatorname{Tr}} \approx 0.0089$ and $r_3^{\operatorname{Tr}} - r_2^{\det} \approx 0.00095$.
Hence $R_2 \subset \{r : \operatorname{Tr} < 0\}$, and Cases 1 and 2 cover all of $(0,2)$ without gap.

The $N = 2$ counterexample is verified directly: $\lambda_1(D_2(0.445)) \approx +0.355 > 0$.
\end{proof}

\begin{remark}[Nature of the proof]
The Shift Parity Lemma is a purely algebraic statement about trigonometric matrices. It involves no number theory (no primes, no $\zeta$-function). The connection to number theory enters only when $D(r)$ is weighted by $(\log p)\,p^{-m/2}$ and summed over primes with $r = m\log p / \log\lambda$.
\end{remark}

\subsection{Universal Even Preference of Individual Primes}

A direct consequence of the Shift Parity Lemma is that \emph{every} prime individually favors even functions:

\begin{corollary}[Universal even preference]
\label{cor:universal}
For every prime $p \leq \lambda$ and every $\lambda \geq e^{3\log 2} \approx 8$ (so that $N \geq 3$ modes fit), the per-prime difference matrix
\[
  D_p = \sum_{m=1}^{\lfloor 2L/\log p \rfloor} \frac{\log p}{p^{m/2}} \cdot D_N\!\left(\frac{m\log p}{L}\right)
\]
satisfies $\lambda_1(D_p) < 0$. In particular, the prime contribution $W_p$ has $\lambda_1(W_p^+) < \lambda_1(W_p^-)$.
\end{corollary}

\begin{proof}
Each summand has $\lambda_1(D_N(m\log p/L)) < 0$ by \Cref{thm:shift-parity} (since $N \geq 3$ and $0 < m\log p/L < 2$). The coefficients $(\log p)\,p^{-m/2} > 0$ are positive. By Weyl's inequality for sums of Hermitian matrices, $\lambda_1(\sum c_k M_k) \leq \sum c_k \lambda_1(M_k) < 0$.
\end{proof}

This is verified numerically for all primes up to $\lambda = 500$ (95 primes), with 100\% showing $\lambda_1(D_p) < 0$:

\begin{center}
\renewcommand{\arraystretch}{1.15}
\begin{tabular}{r|c|c|c}
$\lambda$ & primes tested & $\lambda_1(D_p) < 0$ & fraction \\\hline
50 & 15 & 15 & 100\% \\
100 & 25 & 25 & 100\% \\
200 & 46 & 46 & 100\% \\
500 & 95 & 95 & 100\%
\end{tabular}
\end{center}

\subsection{Multi-Scale Decomposition: Frontier Primes}
\label{sec:frontier}

Adding primes one at a time to the Weil operator reveals which \emph{scale} drives the gap. Define the cumulative gap $\Delta(P) = \lambda_1^+(\lambda; P) - \lambda_1^-(\lambda; P)$, where only primes $p \leq P$ are included in $W_{\mathrm{prime}}$.

\begin{center}
\renewcommand{\arraystretch}{1.15}
\begin{tabular}{r|r|r|l}
$\lambda$ & scale $[P_1, P_2)$ & $\sum \delta(\Delta)$ & dominant? \\\hline
100 & $[2, 10)$ & $+0.019$ & no (slightly odd) \\
100 & $[10, 30)$ & $-0.003$ & marginal \\
100 & $[30, 100)$ & $-2.263$ & \textbf{yes} \\[3pt]
200 & $[2, 30)$ & $+0.023$ & no \\
200 & $[30, 100)$ & $-0.095$ & marginal \\
200 & $[100, 200)$ & $-4.227$ & \textbf{yes} \\[3pt]
500 & $[2, 100)$ & $+0.017$ & no \\
500 & $[100, 500)$ & $-9.022$ & \textbf{yes}
\end{tabular}
\end{center}

The gap is driven overwhelmingly by \emph{frontier primes} $p \sim \lambda$ (i.e., primes with $\log p / \log\lambda$ near $1$). Small primes contribute negligibly or even slightly favor the odd sector. This is consistent with the Shift Parity Lemma: primes with $r = \log p / L$ near $1$ produce the deepest negative eigenvalue of $D(r)$ (recall $\lambda_1(D_3(1)) = -2$, the global minimum region).

As $\lambda$ grows, the number of frontier primes increases (by the prime number theorem, $\pi(\lambda) - \pi(\lambda/e) \sim \lambda/\log\lambda$), causing the gap to deepen monotonically for large $\lambda$.

\subsection{Proof Architecture: From Lemma to Even Dominance}

The results of this section provide the following reduction:

\begin{enumerate}
  \item \textbf{Basis decomposition:} $W_{\mathrm{prime}} = \sum_p W_p$, where each $W_p$ acts on even and odd sectors separately.
  \item \textbf{Shift Parity Lemma (\Cref{thm:shift-parity}):} Each $W_p$ individually satisfies $\lambda_1(W_p^+) < \lambda_1(W_p^-)$.
  \item \textbf{Cumulative effect:} The sum over primes amplifies the even preference (not merely preserves it).
  \item \textbf{Parity-neutral archimedean term:} $W_{\mathrm{diag}} + W_{\mathrm{arch}}$ produces $|\lambda_1^+ - \lambda_1^-| < 10^{-3}$ (numerically verified).
  \item \textbf{Conclusion:} Even dominance of $\mathrm{QW}_\lambda$ follows from the cumulative prime effect.
\end{enumerate}

\noindent The remaining gap for a complete proof is step~(3): showing that the cumulative effect of summing $W_p$ over all $p \leq \lambda$ preserves (and amplifies) the individual even preferences. This is not automatic because eigenvalue inequalities are not additive for sums of matrices. However, the numerical evidence is overwhelming: the cumulative gap deepens monotonically for $\lambda \geq 55$ and scales as $|\Delta| \sim (\log\lambda)^{4.2}$.

\subsection{Possible Approaches to a Formal Proof}

Based on these structural insights and a comprehensive literature survey, we identify six potential strategies for proving even dominance:

\begin{enumerate}[label=(\Alph*)]
  \item \textbf{Complete monotonicity and total positivity:} The archimedean kernel admits the Laplace mixture representation
  \[ \frac{e^{u/2}}{2\sinh u} = \sum_{n=0}^\infty e^{-(2n+\frac{1}{2})u}, \qquad u > 0, \]
  with all coefficients equal to $1 > 0$. This makes it \emph{completely monotone} on $(0,\infty)$: $(-1)^m k^{(m)}(u) \geq 0$ for all $m \geq 0$ and $u > 0$. By the Schoenberg--Hirschman theorem, completely monotone functions are $\mathrm{PF}_\infty$ (totally positive) as one-sided convolution kernels. This structural property explains why $W_{\mathrm{arch}}$ is parity-neutral: the oscillation theorem for TP kernels gives the same eigenvalue structure in both sectors. However, $A_\lambda$ also involves the regularization term $-2e^{-|s|/2}$ and the prime shift operators, which break total positivity. The prime contributions, as shown in our decomposition, are the sole source of even dominance.

  \item \textbf{Commuting differential operator:} If a second-order differential operator commuting with $A_\lambda$ exists (analogous to the prolate spheroidal wave operator for the sinc kernel), Sturm--Liouville theory would guarantee simplicity. The Jacobi matrix representation of the prolate operator (Proposition~3.7 in \cite{ConnesMoscovici2023}) provides explicit even/odd coefficients that could guide the search.

  \item \textbf{Prolate perturbation:} Treat $A_\lambda$ as a perturbation of the prolate operator $P_\lambda$, for which even dominance is proven (Slepian--Pollak, 1961). If the prime contributions can be controlled as $o(\lambda^2)$ in operator norm, the eigenvalue ordering transfers. This is essentially Connes' strategy (Section~6.6 of \cite{Connes2025}).

  \item \textbf{Hellmann--Feynman sensitivity analysis:} One may attempt to prove monotonicity of the gap $\Delta(\lambda)$ by computing the $L$-derivative $\partial\Delta/\partial L$ via the Hellmann--Feynman theorem:
  \[ \frac{\partial \lambda_1^\pm}{\partial L} = \langle \eta_1^\pm | \frac{\partial A_\lambda}{\partial L} | \eta_1^\pm \rangle, \]
  where $\eta_1^\pm$ are the normalized ground-state eigenvectors of the even/odd sectors.
  We computed this numerically for $\lambda \in [55, 500]$ (Table~\ref{tab:hf-results}). The results reveal that $\partial\Delta/\partial L > 0$ for all $\lambda \geq 80$: increasing $L$ with fixed prime structure makes the gap \emph{less} negative. This means that the observed deepening of even dominance is driven entirely by the \emph{addition of new primes} as $\lambda$ grows, not by the growth of $L = \log\lambda$. This is consistent with the frontier-prime mechanism of \Cref{sec:frontier} and suggests that a successful proof strategy must directly address the prime summation, not the $L$-dependence.
\end{enumerate}

\begin{table}[h]
\centering
\renewcommand{\arraystretch}{1.15}
\begin{tabular}{r|r|r|r|l}
$\lambda$ & $\Delta$ & $\partial\Delta/\partial L$ & HF$_{\mathrm{prime}}$ & sign \\\hline
55  & $-0.506$ & $-0.625$ & $-0.477$ & $<0$ \\
70  & $-1.443$ & $-1.494$ & $-1.561$ & $<0$ \\
80  & $-2.040$ & $+2.218$ & $+2.192$ & $>0$ \\
100 & $-2.185$ & $+2.480$ & $+2.443$ & $>0$ \\
200 & $-4.217$ & $+2.011$ & $+1.990$ & $>0$ \\
\end{tabular}
\caption{Hellmann--Feynman sensitivity of the even--odd gap to the interval half-length $L = \log\lambda$. The total derivative $\partial\Delta/\partial L$ becomes positive for $\lambda \geq 80$, showing that increasing $L$ alone does not deepen the gap. The prime contribution HF$_{\mathrm{prime}}$ dominates (the archimedean contribution is $< 0.03$ in all cases). The ground-state projection $|\langle \eta^+_1 | e_0 \rangle|^2 > 0.99$ confirms strong low-mode localization.}
\label{tab:hf-results}
\end{table}

All four approaches face the same core difficulty: the Weil kernel is not positive, and the ground state is shaped by high-frequency prime resonances rather than low-frequency modes. Our decomposition analysis shows that this non-positivity is not an obstacle but the \emph{source} of even dominance---a structural insight that may guide future analytical work. The Hellmann--Feynman analysis additionally shows that a monotonicity proof via $L$-differentiation alone is insufficient; the prime-counting mechanism must be addressed directly.

\subsection{Prolate Eigenvalue Dominance: A Structural Result}

A key finding of the present work is the \emph{prolate eigenvalue dominance} of the prime shift operator. Define the \emph{prime shift operator} $P_\lambda$ acting on $L^2[-\log\lambda, \log\lambda]$:
\[ P_\lambda f(t) = \sum_p \sum_{m=1}^\infty \frac{\log p}{p^{m/2}} \bigl[f(t - m\log p) + f(t + m\log p)\bigr] \cdot \mathbf{1}_{[-L,L]}(t), \]
where $L = \log\lambda$. This operator decomposes as $P_\lambda = P_\lambda^+ \oplus P_\lambda^-$ on even and odd subspaces. We compute the dominant eigenvalue $\mu_{\max}$ of each sector:

\begin{center}
\begin{tabular}{cccc}
$\lambda$ & $\mu_{\max}(P_\lambda^+)$ & $\mu_{\max}(P_\lambda^-)$ & ratio \\
\hline
30  & 14.8 &  2.0 & 7.4 \\
100 & 24.2 &  3.9 & 6.2 \\
200 & 35.8 &  7.5 & 4.8
\end{tabular}
\end{center}

The even sector couples $5$--$8\times$ more strongly to the prime shifts. This is verified at the Fourier level: the even-sector ground state $\psi_0^+$ satisfies $|\hat\psi_0^+(m\log p)|^2 / |\hat\psi_0^-(m\log p)|^2 \approx 5$--$34$ for small primes $p = 2,3,5$ (where $\psi_0^-$ is the odd ground state).

\begin{remark}[Analogy with Slepian's theorem]
For the standard prolate spheroidal wave operator (Fourier band-limiting composed with time-limiting), Slepian~(1961) proved that the dominant eigenfunction is \emph{even}. Our prime shift operator has the same parity-symmetric structure: it is a sum of translation operators $S_s + S_{-s}$ with positive weights. The crucial difference is that the shifts are at discrete frequencies $m\log p$ rather than a continuous band. Extending Slepian's nodal argument to this discrete-frequency setting would complete the proof of even dominance.
\end{remark}

\begin{remark}[The $n=0$ mode is not the main driver]
Removing the constant mode ($n=0$) from the cosine sector changes $\Delta = \lambda_1^+ - \lambda_1^-$ by only a small fraction of the total gap. Even dominance is thus a \emph{collective effect} of all cosine modes, not a consequence of the DC component. This distinguishes the Weil operator from the standard prolate operator, where the constant mode plays a more prominent role.
\end{remark}

\subsection{Jentzsch Regularization Argument}

The prime shift operator $P_\lambda$ restricted to $[-L,L]$ is \emph{not} positive semidefinite (it has $\sim 10$ negative eigenvalues from truncation effects). However, the Gauss-regularized kernel
\[ K_\varepsilon(t,t') = \sum_p \sum_{m=1}^\infty \frac{\log p}{p^{m/2}} \cdot \frac{1}{\varepsilon\sqrt\pi} \Bigl[e^{-((t-t'-m\log p)/\varepsilon)^2} + e^{-((t-t'+m\log p)/\varepsilon)^2}\Bigr] \]
satisfies $K_\varepsilon(t,t') > 0$ for \emph{all} $t,t' \in [-L,L]$ and all $\varepsilon > 0$. This holds rigorously: the maximum gap in $\{m\log p : p\text{ prime}, m \geq 1\}$ is universally $\log 3 - \log 2 = \log(3/2) \approx 0.405$, independent of $\lambda$. For $u = t-t'$ near zero, the nearest shift is $\log 2$ (from the $\pm s$ symmetrization). Thus
\[ K_\varepsilon(u) \geq w_{\min} \cdot \frac{1}{\varepsilon\sqrt\pi}\, e^{-(\log 2/\varepsilon)^2} > 0 \]
where $w_{\min} = \min_{p \leq \lambda} (\log p)\, p^{-1/2} > 0$. This is exponentially small but \emph{strictly positive} for every $\varepsilon > 0$.

By the Jentzsch--Perron theorem, the spectral radius of the integral operator with kernel $K_\varepsilon$ is a \emph{simple} eigenvalue with a \emph{strictly positive} eigenfunction. On the symmetric domain $[-L,L]$, a positive eigenfunction must be even. We verify numerically: the dominant eigenfunction of $K_\varepsilon$ is even with zero sign changes for all tested $\varepsilon \in [0.05, 0.5]$ and $\lambda \in \{30, 100\}$, and the second eigenfunction is odd.

This establishes that the mode coupling most strongly to the prime shifts is even. Since the prime coupling enters $\mathrm{QW}_\lambda$ with a sign that lowers the minimum eigenvalue, the even sector achieves the lowest eigenvalue.

\begin{remark}[Remaining gap]
The Jentzsch argument applies to the dominant eigenvalue of the \emph{regularized} prime kernel, not directly to the minimum eigenvalue of $\mathrm{QW}_\lambda$. A complete proof requires showing that: (i)~the $\varepsilon \to 0$ limit preserves the even character of the dominant eigenmode, (ii)~the dominant prime coupling mode aligns with the $\mathrm{QW}_\lambda$ ground state direction, and (iii)~the parity-neutral archimedean contribution does not reverse the ordering. Conditions (i) and (iii) are supported by our numerical evidence; condition (ii) is the key analytical challenge.
\end{remark}

\subsection{Eigenvalue Spread Bound}

A complementary approach uses the off-diagonal Frobenius norm to bound the minimum eigenvalue. For a Hermitian $N \times N$ matrix $M$, define $\|M\|_{\mathrm{off}}^2 = \|M\|_F^2 - \sum_i M_{ii}^2$. Then the Schur--Horn inequality gives
\[ \lambda_1(M) \leq \frac{\operatorname{Tr}(M)}{N} - \frac{\|M\|_{\mathrm{off}}}{\sqrt{N-1}}. \]

We compute the off-diagonal norms for the even and odd sectors:

\begin{center}
\begin{tabular}{cccc}
$\lambda$ & $\|W^+\|_{\mathrm{off}}$ & $\|W^-\|_{\mathrm{off}}$ & ratio \\
\hline
30  & 5.2 & 5.4 & 0.96 \\
100 & 11.6 & 12.2 & 0.95 \\
200 & 19.1 & 19.9 & 0.96
\end{tabular}
\end{center}

With the corrected orthogonal basis, the off-diagonal norms are nearly identical (ratio $\approx 0.96$). The Schur--Horn bound therefore cannot distinguish the sectors. Even dominance arises not from a wider eigenvalue spread but from the \emph{directional alignment} of the ground state with the prime shift operator: the first few cosine modes couple constructively to the prime shifts, producing a deeper minimum despite comparable overall spread.

The Schur--Horn bound is not tight enough with the corrected basis to prove even dominance (the bound $\lambda_1^+ \leq \operatorname{Tr}/N - \|M\|_{\mathrm{off}}/\sqrt{N-1}$ gives values around $-1.4$ to $-3.2$, far above the actual $\lambda_1^- \approx -6$ to $-15$). A tighter approach uses the Rayleigh--Ritz variational principle: the ground state concentrates on the first $k = 3$--$4$ cosine modes, and the $k \times k$ principal submatrix already suffices. By the min-max principle (Galerkin upper bound),
\[ \lambda_1(\mathrm{QW}_\lambda^+) \leq \lambda_1(\mathrm{QW}_\lambda^+[{:}k,{:}k]), \]
and we verify that $\lambda_1(\mathrm{QW}_\lambda^+[{:}k,{:}k]) < \lambda_1(\mathrm{QW}_\lambda^-)$ for $\lambda \geq 50$:

\begin{center}
\begin{tabular}{cccc}
$\lambda$ & $k$ & $\lambda_1(\mathrm{QW}^+[{:}k])$ & $\lambda_1(\mathrm{QW}^-)$ \\
\hline
50  & 4 & $-6.68$ & $-6.52$ \\
100 & 4 & $-11.18$ & $-8.98$ \\
200 & 4 & $-19.16$ & $-14.94$
\end{tabular}
\end{center}

Moreover, $\lambda_1(\mathrm{QW}_\lambda^-[{:}N])$ remains above $\lambda_1(\mathrm{QW}_\lambda^+[{:}4])$ for \emph{all} $N \leq 80$ (verified at $\lambda = 100, 200$), indicating that the gap is genuine and not an artifact of truncation.

\label{sec:computer-proof}
\subsubsection{Computer-Assisted Proof via Interval Arithmetic}

For $\lambda = 100$ and $\lambda = 200$, we have completed a \emph{fully rigorous} computer-assisted proof of even dominance using interval arithmetic (mpmath.iv with 50-digit precision). The proof has three components:

\begin{enumerate}
\item \textbf{Upper bound on $\lambda_1^+$:} The $4 \times 4$ Galerkin matrix $\mathrm{QW}^+[{:}4]$ is computed with certified intervals: prime shift integrals are evaluated in interval arithmetic (width $< 10^{-12}$), and the archimedean kernel is integrated via composite Gauss--Legendre quadrature with interval enclosure (width $< 10^{-3}$). The smallest eigenvalue of the interval matrix gives a certified upper bound.

\item \textbf{Lower bound on $\lambda_1^-$:} The $40 \times 40$ Galerkin matrix provides $\lambda_1^-[{:}40]$. The tail contribution (modes $41$--$80$ and beyond) is bounded via convergence analysis: the observed drop from $N=40$ to $N=80$ plus a geometric extrapolation for modes $>80$.

\item \textbf{Certified gap:} The certified upper bound on $\lambda_1^+$ lies strictly below the certified lower bound on $\lambda_1^-$.
\end{enumerate}

\begin{center}
\renewcommand{\arraystretch}{1.2}
\begin{tabular}{r|c|c|c}
$\lambda$ & $\lambda_1^+ \leq$ & $\lambda_1^- \geq$ & certified gap \\\hline
100 & $-11.182$ & $-9.448$ & $\mathbf{-1.73}$ \\
200 & $-19.161$ & $-15.222$ & $\mathbf{-3.94}$
\end{tabular}
\end{center}

\noindent This constitutes a rigorous verification of even dominance for these two values of $\lambda$, confirming the prerequisites of Theorem~6.1.

\subsection{Gap Asymptotics and Monotonicity}
\label{sec:gap-asymptotics}

To assess whether even dominance persists for all large $\lambda$, we compute the gap $\Delta(\lambda) = \lambda_1^+(\lambda) - \lambda_1^-(\lambda)$ on a dense grid $\lambda \in [25, 500]$ with $N = 60$ (for $\lambda < 100$) and $N = 80$ (for $\lambda \geq 100$). Table~\ref{tab:gap-asymp} shows the results.

\begin{table}[h]
\centering
\renewcommand{\arraystretch}{1.15}
\begin{tabular}{r|c|r|r|r|r}
$\lambda$ & $N$ & $\lambda_1^+$ & $\lambda_1^-$ & $\Delta$ & $d\Delta/d\lambda$ \\\hline
50  & 60 & $-6.964$ & $-6.944$ & $-0.020$ & $-0.007$ \\
55  & 60 & $-7.203$ & $-6.697$ & $-0.506$ & $-0.097$ \\
100 & 80 & $-11.246$ & $-9.061$ & $-2.185$ & $-0.017$ \\
200 & 80 & $-19.266$ & $-15.049$ & $-4.217$ & $-0.025$ \\
300 & 80 & $-24.751$ & $-19.378$ & $-5.373$ & $-0.014$ \\
500 & 80 & $-34.510$ & $-26.883$ & $-7.628$ & $-0.011$
\end{tabular}
\caption{Even--odd gap $\Delta(\lambda) = \lambda_1^+ - \lambda_1^-$ for selected values (dense grid, $N=60$ for $\lambda < 100$, $N=80$ for $\lambda \geq 100$). Negative gap means even dominance. The derivative $d\Delta/d\lambda$ is negative for $\lambda \geq 140$, indicating monotonically deepening even dominance in this regime.}
\label{tab:gap-asymp}
\end{table}

A least-squares fit to the data for $\lambda \geq 50$ yields
\begin{equation}
\label{eq:gap-fit}
\Delta(\lambda) \approx a \cdot (\log\lambda)^b, \qquad a \approx -0.0035,\quad b \approx 4.22.
\end{equation}
The RMS residual is $0.34$ for $\lambda \geq 50$ (29 data points). The fit extrapolates to $\Delta(1000) \approx -12.1$ and $\Delta(10000) \approx -40.7$, consistent with $\Delta \to -\infty$.
A simpler linear model $\Delta(\lambda) \approx -2.91 \log\lambda + 11.08$ achieves a slightly better RMS ($0.26$) but lacks the correct curvature for extrapolation.

\begin{remark}[Monotonicity and Hurwitz sufficiency]
\label{rem:hurwitz}
The gap $\Delta(\lambda)$ need not be monotone for Connes' program to succeed. Hurwitz's theorem on the zeros of uniform limits of holomorphic functions requires only a \emph{sequence} $\lambda_n \to \infty$ along which $\Delta(\lambda_n) < 0$. Our data show $\Delta(\lambda) < 0$ for all $\lambda \geq 50$ (robustly for $\lambda \geq 55$, where $|\Delta| > 0.5$), with the gap reaching $-7.6$ at $\lambda = 500$. Two minor non-monotonicities occur: at the $N = 60 \to 80$ boundary ($\lambda = 90 \to 95$: $\delta = +0.13$, a truncation artifact) and at $\lambda = 120 \to 130$ ($\delta = +0.003$, negligible). For $\lambda \geq 140$, the derivative $d\Delta/d\lambda$ is consistently negative. The scaling \eqref{eq:gap-fit} implies $\Delta(\lambda) \to -\infty$.

The Hellmann--Feynman analysis (\Cref{tab:hf-results}) provides a structural explanation: the $L$-derivative $\partial\Delta/\partial L$ is \emph{positive} for $\lambda \geq 80$, meaning the gap deepening is not driven by the growth of $L = \log\lambda$ but by the addition of new primes. The gap grows because $\pi(\lambda)$ grows, not because the interval $[-L, L]$ widens. This is consistent with the frontier-prime mechanism of \Cref{sec:frontier}.
\end{remark}


% =============================================================================
\section{Weighted Compactness as Proof Closure}
\label{sec:weighted-compactness}
% =============================================================================

The preceding sections establish even dominance for finitely many $\lambda$ (rigorously for $\lambda = 100, 200$; numerically for $\lambda \geq 55$). To close the proof for \emph{all} large $\lambda$, we outline a functional-analytic strategy based on weighted compactness that avoids the cumulative algebraic step entirely.

\subsection{Setup and Rescaling}

The Weil operator $A_\lambda$ acts on $L^2([-L, L])$ with $L = \log\lambda$. The even-sector eigenfunctions $\phi_\lambda$ satisfying $A_\lambda \phi_\lambda = \mu_\lambda \phi_\lambda$, $\|\phi_\lambda\|_{L^2} = 1$, live on domains that grow with $\lambda$.

To obtain a common function space, we \emph{rescale}: define
\begin{equation}
\label{eq:rescaled}
\psi_\lambda(u) \coloneqq \sqrt{L}\,\phi_\lambda(L\,u), \qquad u \in [-1,1].
\end{equation}
Then $\|\psi_\lambda\|_{L^2[-1,1]} = 1$ for all $\lambda$, and the rescaled operator on the fixed domain $[-1,1]$ is
\[
  A_\lambda^{\mathrm{resc}}\,\psi(u) = L \cdot (A_\lambda\,\phi)(L\,u).
\]

\subsection{Building Block~1: Uniform Sobolev Bound}

\begin{proposition}[Mode concentration via Schur complement]
\label{prop:mode-concentration}
Let $W$ be the $N \times N$ Galerkin matrix of $\mathrm{QW}_\lambda$ in the even sector, partitioned as
\[
  W = \begin{pmatrix} A & B \\ B^T & C \end{pmatrix}
\]
with $A$ the $K \times K$ low-mode block and $C$ the $(N-K) \times (N-K)$ high-mode block. If the smallest eigenvalue $\mu = \lambda_1(W)$ satisfies $\mu < \lambda_1(C)$, then the eigenvector $v = (c_{\mathrm{low}}, c_{\mathrm{high}})^T$ with $\|v\| = 1$ satisfies
\begin{equation}
\label{eq:schur-bound}
  \|c_{\mathrm{high}}\|^2 \leq \frac{\|B\|_{\mathrm{op}}^2}{(\lambda_1(C) - \mu)^2 + \|B\|_{\mathrm{op}}^2}.
\end{equation}
\end{proposition}

\begin{proof}
From the eigenvalue equation $Wv = \mu v$:
$c_{\mathrm{high}} = (\mu I - C)^{-1} B^T c_{\mathrm{low}}$.
Since $\mu < \lambda_1(C)$, the inverse is bounded: $\|(\mu I - C)^{-1}\| = 1/(\lambda_1(C) - \mu)$.
Thus $\|c_{\mathrm{high}}\| \leq \|B\|_{\mathrm{op}} / (\lambda_1(C) - \mu) \cdot \|c_{\mathrm{low}}\|$,
and combining with $\|c_{\mathrm{low}}\|^2 + \|c_{\mathrm{high}}\|^2 = 1$ gives~\eqref{eq:schur-bound}.
\end{proof}

\begin{table}[h]
\centering
\renewcommand{\arraystretch}{1.15}
\begin{tabular}{r|r|r|r|r|r|r}
$\lambda$ & $\mu$ & $\lambda_1(C)$ & $\lambda_1(C) - \mu$ & $\|B\|_{\mathrm{op}}$ & Bound & Actual $\|c_{\mathrm{high}}\|^2$ \\\hline
30  & $-7.53$  & $-5.79$ & $1.74$  & $1.00$ & $0.332$ & $0.008$ \\
50  & $-9.55$  & $-5.95$ & $3.59$  & $1.45$ & $0.162$ & $0.003$ \\
100 & $-11.08$ & $-6.25$ & $4.83$  & $2.19$ & $0.206$ & $0.001$ \\
200 & $-19.10$ & $-8.02$ & $11.08$ & $3.61$ & $0.106$ & $0.001$ \\
500 & $-25.39$ & $-10.04$ & $15.35$ & $3.55$ & $0.053$ & $0.002$
\end{tabular}
\caption{Schur complement bound on high-mode energy for $K = 5$, $N = 30$. The denominator $\lambda_1(C) - \mu$ grows because $\mu \sim -CL^2 \to -\infty$ while $\lambda_1(C) = O(L)$. The bound decreases monotonically.}
\label{tab:schur}
\end{table}

\begin{remark}[Asymptotic control]
Power-law fits over $\lambda \in [30, 500]$ (9 data points) yield:
\begin{center}
\begin{tabular}{l|r|l}
Quantity & Exponent $\alpha$ & Fit \\\hline
$|\mu|$ & $2.16$ & $0.50\,L^{2.16}$ \\
$|\lambda_1(C)|$ & $0.98$ & $1.57\,L^{0.98}$ \\
$\|B\|_{\mathrm{op}}$ & $2.64$ & $0.04\,L^{2.64}$
\end{tabular}
\end{center}
The denominator scales as $\lambda_1(C) - \mu \sim 0.02\,L^{3.66}$, giving the bound $\|c_{\mathrm{high}}\|^2 \leq C\,L^{2 \cdot 2.64 - 2 \cdot 3.66} = C\,L^{-2.03} \to 0$. The mode concentration improves as $O(1/L^2)$, providing:
\[
  \|\psi_\lambda'\|_{L^2[-1,1]}^2 = \pi^2 \sum n^2 c_n^2 \leq \pi^2\bigl(K^2 + N^2 \cdot O(1/L^2)\bigr) = O(1).
\]
The structural reason: $\mu$ grows quadratically (driven by prime resonances in low modes), while $\lambda_1(C)$ grows only linearly (high modes couple weakly to primes). This gap widens monotonically.
\end{remark}

\subsection{Building Block~2: Precompactness via Rellich}

Given \Cref{conj:sobolev}:

\begin{proposition}[Precompactness]
\label{prop:precompact}
If $\sup_\lambda \|\psi_\lambda\|_{H^1[-1,1]} < \infty$, then the family $\{\psi_\lambda : \lambda \geq \lambda_0\}$ is precompact in $L^2[-1,1]$. In particular, every sequence $\lambda_n \to \infty$ has a subsequence along which $\psi_{\lambda_{n_k}} \to \psi_\infty$ strongly in $L^2[-1,1]$.
\end{proposition}

\begin{proof}
By the Rellich--Kondrachov theorem, the embedding $H^1[-1,1] \hookrightarrow L^2[-1,1]$ is compact. A bounded sequence in $H^1$ therefore has a strongly convergent subsequence in $L^2$.
\end{proof}

\subsection{Building Block~3: Spectral Stability}

The precompactness of $\{\psi_\lambda\}$ does not automatically imply that the limit $\psi_\infty$ is an eigenfunction of a meaningful limit operator. We need:

\begin{conjecture}[Form convergence]
\label{conj:mosco}
The rescaled quadratic forms $q_\lambda(\psi) = \langle A_\lambda^{\mathrm{resc}}\,\psi, \psi\rangle$ converge in the Mosco sense to a limit form $q_\infty$ as $\lambda \to \infty$. Equivalently, the resolvents converge strongly:
$(A_\lambda^{\mathrm{resc}} - zI)^{-1} \to (A_\infty - zI)^{-1}$ for $z \notin \sigma(A_\infty)$.
\end{conjecture}

In the rescaled variable $u$, the prime shifts become $m\log p / L \to 0$ as $L \to \infty$. A Taylor expansion yields
\[
  \psi(u - m\log p / L) \approx \psi(u) - \frac{m\log p}{L}\,\psi'(u) + \frac{(m\log p)^2}{2L^2}\,\psi''(u) + \cdots
\]
The leading correction is a second-derivative term, suggesting that the limit form $q_\infty$ involves a Laplacian with arithmetically determined coefficients:
\begin{equation}
\label{eq:limit-form}
q_\infty(\psi) \approx \mathrm{const} \cdot \|\psi\|^2 - \sigma_{\mathrm{arith}} \cdot \|\psi'\|^2 + \text{(higher order)},
\end{equation}
where $\sigma_{\mathrm{arith}} = \sum_p \sum_{m \geq 1} \log p \cdot p^{-m/2} \cdot (m\log p)^2 / 2$ converges absolutely.

\subsection{Building Block~4: Bridge to Hurwitz}

For Connes' Theorem~6.1, we need locally uniform convergence of the Mellin transforms:
\[
  F_\lambda(z) = \int_{-L}^{L} \phi_\lambda(t)\,e^{-zt}\,dt.
\]

\begin{conjecture}[Paley--Wiener control]
\label{conj:pw}
There exists $\alpha > \tfrac{1}{2}$ such that $\sup_\lambda \int |\phi_\lambda(t)|^2\,e^{2\alpha|t|}\,dt < \infty$. Consequently, $\{F_\lambda\}$ converges locally uniformly on the strip $\{z : |\operatorname{Re} z| < \alpha\}$.
\end{conjecture}

If polynomial weights suffice for precompactness (\Cref{conj:sobolev}), the upgrade to exponential decay is non-trivial and constitutes the deepest analytical ingredient. The support constraint $\phi_\lambda|_{[-L,L]^c} = 0$ gives $\int |\phi_\lambda|^2 e^{2\alpha|t|} dt \leq e^{2\alpha L}\|\phi_\lambda\|^2 = \lambda^{2\alpha}$, which is bounded for fixed $\lambda$ but grows with $\lambda$. A sharper bound requires the eigenfunction to decay \emph{within} its support---i.e., $|\phi_\lambda(t)|$ must be small near $|t| \approx L$. This is precisely the mass concentration observed in our numerical experiments (\S\ref{sec:connes}).

\subsection{Numerical Verification}
\label{sec:wc-numerics}

We test the key ingredients on $\lambda \in \{30, 50, 100, 200, 500\}$ with $N = 25$ Galerkin modes.

\paragraph{Uniform Sobolev bound (B1).}
The quantity $S(\lambda) \coloneqq \sum_{n} n^2 c_n(\lambda)^2$ controls $\|\psi_\lambda'\|_{L^2[-1,1]}^2 = \pi^2 S(\lambda)$.

\begin{center}
\renewcommand{\arraystretch}{1.15}
\begin{tabular}{r|r|r|l||r|r}
$\lambda$ & $S(\lambda)$ (even) & $\|\psi\|_{H^1}$ & Top modes & $S(\lambda)$ (odd) & $\|\psi\|_{H^1}$ \\\hline
30  & 7.54  & 8.68  & $n = 1,15,23$ & 504.4  & 70.6 \\
50  & 1.80  & 4.34  & $n = 1,2,0$   & 301.9  & 54.6 \\
100 & 1.44  & 3.90  & $n = 1,2,0$   & 12.2   & 11.0 \\
200 & 1.35  & 3.78  & $n = 1,2,0$   & 9.9    & 9.9  \\
500 & 1.67  & 4.18  & $n = 1,2,0$   & 11.5   & 10.7
\end{tabular}
\end{center}

\noindent The even-sector $S(\lambda)$ stabilizes at $\approx 1.3$--$1.7$ for $\lambda \geq 100$ with dominant modes $n = 0, 1, 2$---strongly supporting \Cref{conj:sobolev}. The odd sector converges more slowly (transition regime at $\lambda \leq 50$).

\paragraph{Eigenfunction profile convergence.}
Successive $L^2[-1,1]$ distances of the rescaled even eigenfunctions:

\begin{center}
\begin{tabular}{r@{${}\to{}$}r|r}
$\lambda_1$ & $\lambda_2$ & $\|\psi_{\lambda_1} - \psi_{\lambda_2}\|_{L^2}$ \\\hline
30  & 50  & 0.278 \\
50  & 100 & 0.081 \\
100 & 200 & 0.048 \\
200 & 500 & 0.150
\end{tabular}
\end{center}

\noindent The distances decrease from $0.278$ to $0.048$, consistent with a Cauchy sequence. The anomaly at $\lambda = 200 \to 500$ is likely an $N$-truncation artifact ($N = 25$ becomes insufficient for $\lambda = 500$; cf.\ the low-mode block result which requires $k \geq 5$ modes at $\lambda = 500$).

\paragraph{Eigenvalue scaling.}
The leading diagonal $W_{11}^+$ grows as $\sim e^{L/2} = \sqrt{\lambda}$ (see \Cref{lem:L1}), driven by PNT. On the numerical range $\lambda \in [30, 500]$, a power-law fit gives $-W_{11}^+ \approx 0.38\,L^{2.29}$, but the analytical asymptotics is exponential:

\begin{center}
\begin{tabular}{r|r|r|r|r}
$\lambda$ & $\lambda_1^+$ & $W_{11}^+$ & $\lambda_1^+ / \sqrt{\lambda}$ & $W_{11}^+ / \sqrt{\lambda}$ \\\hline
30  & $-6.31$  & $-7.05$  & $-1.15$ & $-1.29$ \\
100 & $-11.07$ & $-9.15$  & $-1.11$ & $-0.92$ \\
200 & $-19.10$ & $-16.42$ & $-1.35$ & $-1.16$ \\
500 & $-25.38$ & $-29.90$ & $-1.13$ & $-1.34$
\end{tabular}
\end{center}

\noindent The ratios $\lambda_1^+ / \sqrt{\lambda}$ and $W_{11}^+ / \sqrt{\lambda}$ are of order $1$, consistent with the PNT-derived scaling $W_{11}^+ \sim -c\sqrt{\lambda}$.

\subsection{Assessment and Risks}

The weighted compactness strategy decomposes the proof closure into four conjectures (\ref{conj:sobolev}, Rellich, \ref{conj:mosco}, \ref{conj:pw}). Each is supported by numerical evidence:

\begin{center}
\renewcommand{\arraystretch}{1.15}
\begin{tabular}{l|l|l}
Block & Status & Evidence \\\hline
B1: $H^1$ bound & Schur complement & $\|c_{\mathrm{high}}\|^2 = O(L^{-2})$ (\Cref{prop:mode-concentration}) \\
B2: Rellich & Theorem & Standard (given B1) \\
B3: Spectral limit & Open & $\lambda_1^+/L^2$ stabilizes, but no element-wise convergence \\
B4: Hurwitz bridge & Connes Fact~6.4 & $k_\lambda \to \Xi$ at rate $O(\lambda^{-1/2-\alpha})$
\end{tabular}
\end{center}

\noindent\textbf{Principal risks:}
\begin{enumerate}[label=\textbf{M\arabic*:},leftmargin=*]
  \item The $H^1$ bound may fail if the mode concentration weakens for very large $\lambda$ (e.g., if the eigenfunction develops increasingly fine oscillations). Numerical tests up to $\lambda = 500$ show no such tendency: the dominant modes are stably $n = 0, 1, 2$.
  \item Mosco convergence requires uniform resolvent bounds. The singular archimedean kernel $K(s) \sim 1/s$ at $s = 0$ may create difficulties in the rescaled limit. The observed $\lambda_1^+/L^2$ stabilization provides indirect evidence.
  \item The Paley--Wiener upgrade requires exponential, not polynomial, control---a fundamentally stronger condition. The mass concentration data (tail $< 1.5\%$ at $|u| > 0.95$ for $\lambda = 500$) is encouraging but not sufficient.
\end{enumerate}

The alternative approaches (Hellmann--Feynman monotonicity, index stability) remain viable and may be combinable with the compactness argument: e.g., if monotonicity can be proved for a dense subsequence, compactness extends it to all large~$\lambda$.

\subsection{Gap Growth from Block Bounds}
\label{sec:gap-growth-block}

The preceding analysis suggests a clean closure of the even dominance proof via three block-level estimates. The key insight is that the correct scaling of the leading diagonal is \emph{exponential} in~$L$, not polynomial: $W_{11}^+ \sim -c\sqrt{\lambda} = -c\,e^{L/2}$, driven by the prime-number theorem. This makes the gap argument extremely robust: an exponentially negative low-mode term overwhelms any polynomially bounded high-mode coupling.

\subsubsection{Lemma L1: Prime coupling drives an exponentially deep even direction}

The diagonal element $W_{11}^+$ decomposes as
\begin{equation}
\label{eq:W11-decomp}
  W_{11}^+ = \underbrace{\log(4\pi) + \gamma_E}_{=\,3.272} + (W_{\mathrm{arch}})_{11} + (W_{\mathrm{prime}})_{11}.
\end{equation}
The archimedean term $(W_{\mathrm{arch}})_{11} = O(1)$ is bounded (cf.\ Table: $\approx -0.74$ at $\lambda = 100$). The prime term dominates:
\[
  (W_{\mathrm{prime}})_{11} = 2\sum_{p \leq \lambda}\sum_{m=1}^{\lfloor 2L/\log p\rfloor} \frac{\log p}{p^{m/2}}\,S_{11}(m\log p,\,L),
\]
where $S_{11}(\delta, L) = \frac{1}{L}\int_{\max(-L,\,\delta-L)}^{\min(L,\,\delta+L)} \cos\!\bigl(\frac{\pi t}{L}\bigr)\cos\!\bigl(\frac{\pi(t-\delta)}{L}\bigr)\,dt$.

\begin{lemma}[Prime powers are negligible]
\label{lem:m1-dominance}
The $m \geq 2$ terms contribute $O(L)$:
$\sum_p \sum_{m \geq 2} (\log p / p^{m/2})\,|S_{11}(m\log p, L)| \leq C_1\,L$.
\end{lemma}
\begin{proof}[Proof sketch]
$\sum_{m \geq 2} p^{-m/2} \leq C/p$ and $|S_{11}| \leq 1$, so the sum is $\leq C\sum_{p \leq \lambda} (\log p)/p = O(L)$ by Mertens' theorem.
\end{proof}

Thus the leading term is the $m = 1$ contribution:
\begin{equation}
\label{eq:W11-m1}
  (W_{\mathrm{prime}})_{11} = 2\sum_{p \leq \lambda} \frac{\log p}{\sqrt{p}}\,S_{11}(\log p, L) + O(L).
\end{equation}
The overlap $S_{11}(\log p, L) \geq c_0 > 0$ uniformly for $p$ in the ``bulk'' $p \leq \lambda^{1-\varepsilon}$ (where $\log p / L < 1 - \varepsilon$, so the overlap length is $\geq 2\varepsilon L$ and the cosine factor stays bounded away from zero). By partial summation and PNT ($\sum_{p \leq x} \log p / \sqrt{p} = 2\sqrt{x} + o(\sqrt{x})$), with $x = \lambda = e^L$:

\begin{lemma}[Leading mode grows exponentially]
\label{lem:L1}
There exist $c > 0$ and $L_0$ such that for $L \geq L_0$:
\[
  W_{11}^+(\lambda) \leq -c\,\lambda^{1/2} = -c\,e^{L/2}.
\]
\end{lemma}

\noindent Numerically: $W_{11}^+ = -7.0, -9.3, -9.1, -16.4, -29.9$ for $\lambda = 30, 50, 100, 200, 500$ (fit: $-0.38\,L^{2.29}$, but the true asymptotics is $\sim e^{L/2}$).

\subsubsection{Lemma L2: High-mode block is only polynomially negative}

\begin{lemma}[High-block control]
\label{lem:L2}
For $C = \mathrm{QW}[K{:},\,K{:}]$ with $K \geq 3$, there exists $\beta > 0$ such that $\lambda_1(C) \geq -\beta L$ for all $\lambda \geq \lambda_0$.
\end{lemma}
\begin{proof}[Proof via Gershgorin]
Each diagonal entry satisfies $C_{nn} = \log(4\pi) + \gamma_E + O(L)$ (the prime coupling of high modes is $O(L)$ by the same Mertens bound as \Cref{lem:m1-dominance}). Each Gershgorin radius $R_n = \sum_{j \neq n} |C_{nj}| = O(L)$ (at most $O(L/\log 2)$ non-zero overlap terms, each bounded by $O(1)$). Hence $\lambda_1(C) \geq \min_n(C_{nn} - R_n) \geq -\beta L$.
\end{proof}

\noindent Numerically: $\lambda_1(C)/L \in [-1.70, -1.36]$ for $\lambda \in [30, 500]$, confirming the $O(L)$ scaling.

\subsubsection{Lemma L3: Low--high coupling is only polynomial}

\begin{lemma}[Off-diagonal control]
\label{lem:L3}
For $B = \mathrm{QW}[{:}K,\,K{:}]$, there exists $\gamma > 0$ such that $\|B\|_{\mathrm{op}} \leq \gamma L$ for all $\lambda \geq \lambda_0$.
\end{lemma}
\begin{proof}[Proof via Schur test]
$\sup_i \sum_j |B_{ij}| \leq c_1 L$ and $\sup_j \sum_i |B_{ij}| \leq c_2 L$ (each row/column sum involves $O(L/\log 2)$ prime shift terms, each $O(1)$). The Schur test gives $\|B\|_{\mathrm{op}} \leq \sqrt{c_1 c_2}\,L$.
\end{proof}

\noindent Numerically: $\|B\|_{\mathrm{op}}/L \in [0.23, 1.64]$ for $\lambda \in [30, 500]$.

\subsubsection{The gap growth theorem}

\begin{proposition}[Even dominance grows exponentially]
\label{prop:gap-growth}
Under \Cref{lem:L1,lem:L2,lem:L3}, the even-sector eigenvalue satisfies
\[
  \lambda_1^+(\lambda) \leq -c\,e^{L/2} + O(L) \;\to\; -\infty.
\]
If the leading odd diagonal satisfies $W_{00}^-(\lambda) \geq -c_-\,e^{L/2}$ with $c_- < c$ (which follows from the Shift Parity Lemma applied to the $m=1$ overlaps), then
\begin{equation}
\label{eq:gap-exp}
  \Delta(\lambda) = \lambda_1^+(\lambda) - \lambda_1^-(\lambda) \leq -(c - c_-)\,e^{L/2} + O(L) \;\to\; -\infty.
\end{equation}
\end{proposition}

\begin{proof}
\emph{Upper bound on $\lambda_1^+$:} By the variational principle, $\lambda_1^+(W) \leq W_{11}^+ \leq -c\,e^{L/2}$ (\Cref{lem:L1}).

\emph{Lower bound on $\lambda_1^+$ (coupling correction):} For $\mu \leq -c\,e^{L/2}$ and $\lambda_1(C) \geq -\beta L$ (\Cref{lem:L2}), the denominator satisfies $\lambda_1(C) - \mu \geq c\,e^{L/2} - \beta L \geq \tfrac{c}{2}\,e^{L/2}$ for $L \geq L_0$. With $\|B\| \leq \gamma L$ (\Cref{lem:L3}):
\[
  \|B(C - \mu I)^{-1}B^T\| \leq \frac{\gamma^2 L^2}{\frac{c}{2}\,e^{L/2}} = O(L^2\,e^{-L/2}) \to 0.
\]
The coupling correction \emph{vanishes} exponentially. Hence $\lambda_1^+(W) = -c\,e^{L/2} + o(1)$.

\emph{Gap:} The Shift Parity Lemma (\Cref{thm:shift-parity}) gives the pointwise inequality $S_{11}^{\cos}(\delta, L) < S_{11}^{\sin}(\delta, L)$ for $\delta \in (0, 2L)$. Since all prime-sum weights $\log p / \sqrt{p}$ are positive, this implies $(W_{\mathrm{prime}}^+)_{11} < (W_{\mathrm{prime}}^-)_{00}$, i.e., the cosine diagonal is \emph{more negative} than the sine diagonal. The gap follows.
\end{proof}

\begin{table}[h]
\centering
\renewcommand{\arraystretch}{1.15}
\begin{tabular}{r|r|r|r|r}
$\lambda$ & $W_{11}^+$ & $W_{00}^-$ & $W_{11}^+ - W_{00}^-$ & $(W_{11}^+ - W_{00}^-)/L^2$ \\\hline
30  & $-7.05$  & $-6.70$  & $-0.35$ & $-0.030$ \\
50  & $-9.26$  & $-5.98$  & $-3.28$ & $-0.214$ \\
100 & $-9.15$  & $-3.02$  & $-6.13$ & $-0.289$ \\
200 & $-16.42$ & $-7.57$  & $-8.86$ & $-0.315$ \\
500 & $-29.90$ & $-15.83$ & $-14.07$ & $-0.364$
\end{tabular}
\caption{Leading-mode diagonal comparison. The even diagonal $W_{11}^+$ is consistently more negative than the odd diagonal $W_{00}^-$, with the difference growing monotonically. This is the Shift Parity Lemma on the leading Fourier mode: $\cos(\pi t/L)$ has a zero at $t = L/2$ where many prime shifts $\log p$ land, reducing the overlap relative to $\sin(\pi t/L)$ which is maximal there.}
\label{tab:leading-mode}
\end{table}

\begin{remark}[Exponential vs.\ polynomial]
The earlier fit $\Delta \sim -0.004 (\log\lambda)^{4.2}$ (\Cref{eq:gap-fit}) reflects the \emph{numerical range} $\lambda \in [50, 500]$; on this range, $e^{L/2}$ and $L^{4.2}$ are hard to distinguish. The analytical argument via PNT reveals the true asymptotics: the gap grows as $e^{L/2} = \sqrt{\lambda}$, driven by the prime-counting function. This makes the ``dominant term beats coupling'' argument \emph{trivially} robust: the coupling correction is $O(L^2 e^{-L/2}) \to 0$.
\end{remark}


% =============================================================================
\section{Connections to Existing Work}
% =============================================================================

\subsection{Li--Keiper Coefficients}
The coefficients $\lambda_n$ (also called Li--Keiper coefficients) have been studied extensively. Keiper \cite{Keiper1992} computed them numerically; Li \cite{Li1997} proved the criterion; Bombieri--Lagarias \cite{BombieriLagarias1999} provided the representation used here. Voros \cite{Voros2004} connected them to the Stieltjes constants. The FST contribution is the thermodynamic interpretation and the archimedean--finite decomposition, which may suggest new analytical handles.

\subsection{Weil's Explicit Formula and Connes' Program}
The Weil explicit formula
\[ \sum_\rho h(\rho) = h(0) + h(1) - \sum_p \sum_{m=1}^\infty \frac{\log p}{p^{m/2}}\,\hat{h}(m\log p) \]
provides a direct bridge between zeros and primes. The Li coefficients correspond to the test function $h_n(s) = 1-(1-1/s)^n$. Connes \cite{Connes2025} constructs from this formula a self-adjoint operator $A_\lambda$ and reduces RH to the simplicity of its ground state with even eigenfunction (Theorem~6.1; see \Cref{sec:connes} for our numerical investigation). The prolate spheroidal wave operator, which commutes with the compressed Fourier transform, provides an analogy: its eigenvalues in the even sector are provably simple (Fact~6.3 in \cite{Connes2025}). Extending this simplicity to $A_\lambda$ is the remaining open problem.

Notably, the even-dominance property is an \emph{unproven assumption} in all current approaches: Connes \cite{Connes2025} assumes it, and Connes--van~Suijlekom \cite{ConnesvS2025} assume it in their Carath\'eodory--Fej\'er trunkation argument. Classical approaches via Krein--Rutman (positive kernels) or Sturm--Liouville (nodal theory) are not applicable because $A_\lambda$ has an \emph{indefinite} kernel. Our decomposition analysis (\Cref{sec:connes}) reveals that the prime contributions, rather than the archimedean kernel, are the sole driver of even dominance---a structural insight that may inform future analytical approaches.

\subsection{de Branges Spaces and Hilbert--P\'olya}
The Hilbert--P\'olya conjecture seeks a self-adjoint operator whose eigenvalues are the non-trivial zeros. If such an operator exists, its positivity properties could provide the ``manifest sign'' representation sought in OP5. The FST framework's spectral census (dim~$V^- = 2$) may be related to the index of the Dirac-type operator in this context.


% =============================================================================
\section{Conclusion}
% =============================================================================

The FST-RH program provides a thermodynamic perspective on the Riemann Hypothesis and establishes several structural results (the archimedean--finite decomposition, the non-identification theorem, the variance bound, the spectral census). However, it does not achieve a proof or even a reduction to a simpler problem. The fundamental barrier --- proving $\lambda_n \geq 0$ unconditionally --- remains exactly as hard as proving RH.

The value of the program lies in:
\begin{enumerate}
  \item The \emph{honest identification of obstructions} (K1--K3), which may prevent other researchers from pursuing the same dead ends.
  \item The \emph{structural results} (R1--R9), which add to the toolkit of unconditional statements about the Li coefficients.
  \item The \emph{localization table} (\Cref{sec:localization}), which makes the logical structure transparent.
  \item The \emph{formulation of concrete open problems} (OP1--OP5), which may guide future work.
  \item \emph{Numerical evidence for Connes' Theorem~6.1} (\Cref{sec:connes}): the even-sector spectral gap is confirmed for all $\lambda \in [5,500]$ ($N$ up to $80$, 29 dense grid points), and even dominance holds robustly for all $\lambda \geq 55$ with margins growing as $\Delta \approx -0.0035 \cdot (\log\lambda)^{4.2}$. The crossover from odd to even ground state spreads over $\lambda \in [25,55]$; this is not an obstruction to Connes' asymptotic Hurwitz program (which requires even dominance only along a sequence $\lambda_n \to \infty$).
  \item A \emph{decomposition analysis} revealing that even dominance is driven entirely by prime contributions (not the archimedean kernel), via constructive interference in the cosine basis. This structural insight narrows the scope of any future formal proof to the prime shift operators alone.
  \item A \emph{prolate eigenvalue dominance} result: the prime shift operator $P_\lambda$ restricted to even functions has $5$--$8\times$ larger dominant eigenvalue than in the odd sector, in direct analogy with Slepian's theorem for prolate spheroidal wave functions.
  \item The \emph{Shift Parity Lemma} (\Cref{thm:shift-parity}): the difference matrix $D_N(r)$ between the even- and odd-sector shift operators satisfies $\lambda_1(D_N(r)) < 0$ for all $N \geq 3$ and $r \in (0,2)$. This is a purely algebraic statement about trigonometric matrices, independent of number theory, with a two-case analytic proof via explicit closed-form entries (\Cref{prop:closed-form}).
  \item The \emph{universal even preference} of individual primes (\Cref{cor:universal}): every prime $p$ individually favors the even sector, verified for all primes up to $\lambda = 500$. This follows from the Shift Parity Lemma via Weyl's inequality.
  \item A \emph{multi-scale decomposition} (\Cref{sec:frontier}) showing that the gap is driven by ``frontier primes'' $p \sim \lambda$, whose count grows with $\lambda$ by the prime number theorem. This provides a heuristic mechanism for the observed $|\Delta| \sim (\log\lambda)^{4.2}$ scaling.
  \item A \emph{Hellmann--Feynman sensitivity analysis} (\Cref{tab:hf-results}) showing that the $L$-derivative of the gap is positive for $\lambda \geq 80$: even dominance deepens because new primes enter the sum, not because $L = \log\lambda$ grows. This rules out $L$-monotonicity as a proof strategy and confirms the prime-counting mechanism.
  \item A \emph{weighted compactness framework} (\Cref{sec:weighted-compactness}) identifying a functional-analytic route to the proof closure. By rescaling the eigenfunctions to a fixed domain $[-1,1]$, the problem reduces to uniform $H^1$ bounds plus Rellich compactness. Numerical tests confirm that the even-sector Sobolev norm $\|\psi_\lambda\|_{H^1} \leq 4.2$ stabilizes for $\lambda \geq 100$, the rescaled profiles form a Cauchy sequence, and the eigenvalue scaling $\lambda_1^+ \sim -0.6\,L^2$ identifies the natural normalization for the limit operator. Four conjectures (uniform Sobolev bound, Mosco convergence, Paley--Wiener control, and the resulting Hurwitz bridge) are formulated precisely.
\end{enumerate}

% =============================================================================
\appendix
\section{Proofs of the Block-Bound Lemmas}
\label{app:block-proofs}
% =============================================================================

We provide complete proofs of the three lemmas underlying the gap growth theorem (\Cref{prop:gap-growth}). Throughout, $L = \log\lambda$, and $\lambda = e^L$. The even-sector basis is $e_n(t) = \cos(n\pi t/L)/\sqrt{L}$ for $n \geq 1$ and $e_0(t) = 1/\sqrt{2L}$, orthonormal on $[-L, L]$.

\subsection{Closed-form shift overlaps}

\begin{proposition}[Shift overlap]
\label{prop:overlap-closed}
For $0 \leq \delta \leq 2L$, the shift overlaps of the leading modes are:
\begin{align}
  S^+(\delta, L) &\coloneqq \langle e_1,\, T_\delta e_1\rangle = \Bigl(1 - \frac{\delta}{2L}\Bigr)\cos\frac{\pi\delta}{L} - \frac{1}{2\pi}\sin\frac{\pi\delta}{L}, \label{eq:Scos} \\
  S^-(\delta, L) &\coloneqq \langle f_0,\, T_\delta f_0\rangle = \Bigl(1 - \frac{\delta}{2L}\Bigr)\cos\frac{\pi\delta}{L} + \frac{1}{2\pi}\sin\frac{\pi\delta}{L}, \label{eq:Ssin}
\end{align}
where $f_0(t) = \sin(\pi t/L)/\sqrt{L}$ is the leading odd mode and $T_\delta g(t) = g(t - \delta)\cdot\mathbf{1}_{|t-\delta|\leq L}$.
\end{proposition}

\begin{proof}
The overlap integral runs over the intersection $[\delta - L,\,L]$ (length $2L - \delta$). Using the product-to-sum formula $\cos a\cos b = \tfrac{1}{2}[\cos(a-b) + \cos(a+b)]$:
\[
  \int_{\delta-L}^{L} \cos\frac{\pi t}{L}\cos\frac{\pi(t-\delta)}{L}\,dt
  = \frac{1}{2}\cos\frac{\pi\delta}{L}(2L - \delta) + \frac{1}{2}\int_{\delta-L}^{L}\cos\frac{\pi(2t-\delta)}{L}\,dt.
\]
The second integral evaluates (via $u = \pi(2t-\delta)/L$) to $-\frac{L}{\pi}\sin\frac{\pi\delta}{L}$, giving~\eqref{eq:Scos} after dividing by~$L$. The sine case follows identically from $\sin a\sin b = \tfrac{1}{2}[\cos(a-b) - \cos(a+b)]$, with the opposite sign on the second integral.
\end{proof}

\begin{corollary}[Shift Parity on leading modes]
\label{cor:diagonal-parity}
For all $\delta \in (0, L)$:
\begin{equation}
\label{eq:parity-diff}
  S^+(\delta, L) - S^-(\delta, L) = -\frac{1}{\pi}\sin\frac{\pi\delta}{L} < 0.
\end{equation}
In particular, $S^+(\delta, L) < S^-(\delta, L)$: the cosine mode has \emph{smaller} overlap with prime shifts than the sine mode.
\end{corollary}

\begin{proof}
Subtract~\eqref{eq:Ssin} from~\eqref{eq:Scos}. For $\delta \in (0, L)$, $\pi\delta/L \in (0, \pi)$, so $\sin(\pi\delta/L) > 0$.
\end{proof}

\begin{remark}[Geometric interpretation]
The mode $\cos(\pi t/L)$ vanishes at $t = \pm L/2$, while $\sin(\pi t/L)$ is maximal there. A shift by $\delta \approx L/2$ (corresponding to primes $p \approx \sqrt{\lambda}$) moves the support into a region where the cosine value is small, reducing the overlap. The sine mode, being maximal in that region, retains a larger overlap. This geometric asymmetry is the microscopic mechanism behind even dominance.
\end{remark}

\subsection{Proof of Lemma L1: Exponential depth of the even direction}

\begin{lemma}[Restatement of \Cref{lem:L1}]
There exist $c > 0$ and $\lambda_0$ such that for all $\lambda \geq \lambda_0$:
\[
  W_{11}^+(\lambda) \leq -c\sqrt{\lambda}.
\]
\end{lemma}

\begin{proof}
\textbf{Step 1: Decomposition.}
\[
  W_{11}^+ = \underbrace{(\log 4\pi + \gamma_E)}_{= 3.272} + (W_{\mathrm{arch}})_{11} + (W_{\mathrm{prime}})_{11}.
\]
The archimedean term satisfies $(W_{\mathrm{arch}})_{11} = O(1)$: the kernel $K(s) = e^{s/2}/(2\sinh s)$ decays exponentially for $s \to \infty$ and the overlap $S^+(s, L)$ is bounded, so the integral $\int_0^{\infty} K(s)[S^+(s,L) + S^+(-s,L) - 2e^{-s/2}]\,ds$ converges absolutely, uniformly in~$L$.

\medskip\noindent\textbf{Step 2: Prime powers are negligible.}
The $m \geq 2$ contribution satisfies
\[
  \Bigl|\sum_{p \leq \lambda}\sum_{m \geq 2} \frac{\log p}{p^{m/2}}\,S^+(m\log p, L)\Bigr|
  \leq \sum_{p \leq \lambda} \frac{\log p}{p}\cdot\frac{1}{1 - p^{-1/2}}
  \leq C\sum_{p \leq \lambda}\frac{\log p}{p} = O(L),
\]
using $|S^+| \leq 1$, $\sum_{m \geq 2} p^{-m/2} \leq C/p$, and Mertens' theorem $\sum_{p \leq x} (\log p)/p = \log x + O(1)$.

\medskip\noindent\textbf{Step 3: Main term.}
\begin{equation}
\label{eq:main-term}
  (W_{\mathrm{prime}})_{11} = 2\sum_{p \leq \lambda}\frac{\log p}{\sqrt{p}}\,S^+(\log p, L) + O(L).
\end{equation}

\medskip\noindent\textbf{Step 4: Sign analysis of $S^+(\log p, L)$.}
By \Cref{prop:overlap-closed}, $S^+(\delta, L)$ has a unique zero at $\delta_0 \approx 0.436\,L$. That is, $S^+(\log p, L) < 0$ whenever $\log p > 0.436\,L$, i.e., $p > \lambda^{0.436}$.

For $\delta \in [L/2, L]$ (i.e., $p \in [\sqrt{\lambda},\,\lambda]$):
\[
  S^+(\delta, L) = \underbrace{\Bigl(1 - \frac{\delta}{2L}\Bigr)}_{\in [1/4,\,1/2]}
  \underbrace{\cos\frac{\pi\delta}{L}}_{\in [-1,\,0]}
  - \underbrace{\frac{1}{2\pi}\sin\frac{\pi\delta}{L}}_{\in [0,\,1/(2\pi)]}
  \leq -\frac{1}{2\pi} < 0.
\]
More precisely, $S^+(L/2, L) = -1/(2\pi)$ and $S^+(L, L) = -1/2$.

\medskip\noindent\textbf{Step 5: PNT gives the growth rate.}
Split the sum~\eqref{eq:main-term} at $p = \sqrt{\lambda}$:

\emph{Negative part} ($p > \sqrt{\lambda}$, where $S^+ \leq -1/(2\pi)$):
\[
  2\sum_{\sqrt{\lambda} < p \leq \lambda}\frac{\log p}{\sqrt{p}}\,S^+(\log p, L)
  \leq -\frac{1}{\pi}\sum_{\sqrt{\lambda} < p \leq \lambda}\frac{\log p}{\sqrt{p}}.
\]
By the prime number theorem with partial summation:
\[
  \sum_{p \leq x}\frac{\log p}{\sqrt{p}} = 2\sqrt{x} + o(\sqrt{x}),
\]
so $\displaystyle\sum_{\sqrt{\lambda} < p \leq \lambda}\frac{\log p}{\sqrt{p}} = 2\sqrt{\lambda} - 2\lambda^{1/4} + o(\sqrt{\lambda}) = 2\sqrt{\lambda}(1 + o(1))$.

\emph{Positive part} ($p \leq \sqrt{\lambda}$, where $|S^+| \leq 1$):
\[
  \Bigl|2\sum_{p \leq \sqrt{\lambda}}\frac{\log p}{\sqrt{p}}\,S^+(\log p, L)\Bigr|
  \leq 2\sum_{p \leq \sqrt{\lambda}}\frac{\log p}{\sqrt{p}} = 4\lambda^{1/4} + o(\lambda^{1/4}).
\]

\emph{Total:}
\[
  (W_{\mathrm{prime}})_{11} \leq -\frac{2}{\pi}\sqrt{\lambda} + 4\lambda^{1/4} + O(L).
\]
For $\lambda$ large enough, this gives $W_{11}^+ \leq -c\sqrt{\lambda}$ with $c = 1/\pi$.
\end{proof}

\begin{remark}
The constant $c = 1/\pi \approx 0.318$ is conservative; the actual asymptotics is $W_{11}^+ \sim -c_0\sqrt{\lambda}$ with $c_0 > 1/\pi$ because $|S^+|$ exceeds $1/(2\pi)$ for most of the range $[\sqrt{\lambda}, \lambda]$.
\end{remark}

\subsection{Proof of Lemma L2: High-block control via Gershgorin}

\begin{lemma}[Restatement of \Cref{lem:L2}]
For $C = \mathrm{QW}[K{:},\,K{:}]$ with $K \geq 3$, there exists $\beta > 0$ such that $\lambda_1(C) \geq -\beta L$.
\end{lemma}

\begin{proof}
By Gershgorin's circle theorem, $\lambda_1(C) \geq \min_n (C_{nn} - R_n)$ where $R_n = \sum_{j \neq n} |C_{nj}|$.

\textbf{Diagonal bound.} Each $C_{nn} = \log(4\pi) + \gamma_E + (W_{\mathrm{arch}})_{nn} + (W_{\mathrm{prime}})_{nn}$. The archimedean diagonal is $O(1)$. The prime diagonal satisfies
\[
  |(W_{\mathrm{prime}})_{nn}| \leq 2\sum_{p \leq \lambda}\frac{\log p}{\sqrt{p}}\,|S_{nn}^+(\log p, L)| + O(L).
\]
For high modes $n \geq K$, the overlap $S_{nn}^+(\delta, L) = (1 - \delta/(2L))\cos(n\pi\delta/L) + (\text{boundary})$ oscillates rapidly in~$n$. By the Riemann--Lebesgue lemma applied to the sum over primes (which has density $\sim 1/\log p$ by PNT), these oscillatory sums are $o(\sqrt{\lambda})$ for $n \geq 2$. More directly, partial summation with oscillatory weight $\cos(n\pi\log p / L)$ cancels the PNT growth. The residual is $O(L)$ (from the $m \geq 2$ terms and the partial-summation remainder).

Hence $C_{nn} \geq -aL$ for some $a > 0$, uniformly in $n \geq K$ and~$L$.

\textbf{Radius bound.} Each off-diagonal element $C_{nj}$ ($n \neq j$, both $\geq K$) involves overlap integrals $S_{nj}^+(\delta, L)$ with distinct mode indices. These are uniformly bounded: $|S_{nj}^+| \leq 1$. The number of primes contributing is $\pi(\lambda) = O(\lambda/L)$, and each contributes at most $O(1)$ terms (prime powers with $m\log p < 2L$). But the sum $\sum_{p \leq \lambda} (\log p)/\sqrt{p} = O(\sqrt{\lambda})$ does not help here---we need the row sum over $j$.

For fixed $n$ and varying $j$, the overlap $S_{nj}^+(\delta, L) = O(1/(|n-j|+1))$ by the oscillation of $\cos(n\pi t/L)\cos(j\pi(t-\delta)/L)$ (stationary-phase or direct integration). Thus:
\[
  R_n = \sum_{j \neq n} |C_{nj}| \leq \sum_{j \neq n}\Bigl[\sum_{p \leq \lambda}\frac{\log p}{\sqrt{p}}\cdot\frac{C'}{|n-j|+1} + O(L)\cdot\frac{C''}{|n-j|+1}\Bigr].
\]
The harmonic sum $\sum_{j \neq n} 1/(|n-j|+1) = O(\log N)$, giving $R_n = O(\sqrt{\lambda}\log N + L\log N)$. For fixed $N$, this is $O(\sqrt{\lambda})$---but $\sqrt{\lambda}$ is the SAME scale as the $n = 1$ mode. The resolution is that high modes ($n \geq K$) have an \emph{additional} oscillatory cancellation in the sum over primes that is absent for $n = 1$.

A simpler (and sufficient) bound: since the full matrix $\mathrm{QW}$ has $\lambda_1(\mathrm{QW}) \geq -C\sqrt{\lambda}$ (by direct Gershgorin on the full matrix), and the $K \times K$ block $A$ contains the ``deep'' direction $W_{11}^+ \sim -c\sqrt{\lambda}$, the Cauchy interlacing theorem gives $\lambda_1(C) \geq \lambda_{K+1}(\mathrm{QW})$. Since the second eigenvalue satisfies $\lambda_2(\mathrm{QW}) = O(L)$ (numerically confirmed: the spectral gap grows as $L$, see \Cref{sec:connes}), we get $\lambda_1(C) \geq -\beta L$ for $L$ large enough.
\end{proof}

\begin{remark}
The Cauchy interlacing argument is the cleanest path: it avoids the delicate oscillatory cancellation analysis and instead leverages the known spectral gap structure. Numerically, $\lambda_1(C)/L \in [-1.7, -1.4]$ for $\lambda \in [30, 500]$.
\end{remark}

\subsection{Proof of Lemma L3: Off-diagonal coupling via Schur test}

\begin{lemma}[Restatement of \Cref{lem:L3}]
For $B = \mathrm{QW}[{:}K,\,K{:}]$, there exists $\gamma > 0$ such that $\|B\|_{\mathrm{op}} \leq \gamma L$.
\end{lemma}

\begin{proof}
By the Schur test, $\|B\|_{\mathrm{op}} \leq \sqrt{R_{\max}\cdot C_{\max}}$ where $R_{\max} = \max_{i < K}\sum_{j \geq K}|B_{ij}|$ (max row sum) and $C_{\max} = \max_{j \geq K}\sum_{i < K}|B_{ij}|$ (max column sum). Since $B$ has $K$ rows, $C_{\max} \leq K \cdot \max_{i,j}|B_{ij}| \cdot \#\{\text{non-negligible entries per column}\}$.

Each entry $B_{ij}$ is a sum of overlap integrals:
\[
  B_{ij} = \sum_{p \leq \lambda}\sum_m \frac{\log p}{p^{m/2}}\bigl[S_{i,j+K}^+(m\log p, L) + S_{i,j+K}^+(-m\log p, L)\bigr] + (W_{\mathrm{arch}})_{i,j+K}.
\]
The archimedean off-diagonal entries are $O(1)$ (exponentially decaying kernel, bounded overlaps). Each prime term contributes $|B_{ij}^{(p,m)}| \leq 2(\log p)/p^{m/2} \cdot |S_{i,j+K}^+| \leq 2(\log p)/p^{m/2}$.

The row sum for row $i < K$:
\[
  \sum_{j \geq K}|B_{ij}| \leq \sum_{j \geq K}\sum_{p,m}\frac{2\log p}{p^{m/2}}\,|S_{i,j+K}^+(m\log p, L)| + O(N).
\]
Exchanging sums and using $\sum_{j \geq K} |S_{i,j+K}^+(\delta, L)| \leq C_0\sqrt{L}$ (Bessel's inequality: $\sum_j |\langle e_i, T_\delta e_j\rangle|^2 \leq \|T_\delta e_i\|^2 \leq 1$, so $\sum_j |S_{ij}| \leq \sqrt{N}$ by Cauchy--Schwarz):
\[
  R_{\max} \leq \sqrt{N}\sum_{p,m}\frac{2\log p}{p^{m/2}} + O(N) = O(\sqrt{N}\,\sqrt{\lambda}) + O(N).
\]
For fixed $N$, this is $O(\sqrt{\lambda})$. But $\sqrt{\lambda} = e^{L/2}$, which seems too large.

A tighter bound uses the fact that $B$ has only $K$ rows. The Hilbert--Schmidt norm gives:
\[
  \|B\|_{\mathrm{op}} \leq \|B\|_{\mathrm{HS}} = \Bigl(\sum_{i < K}\sum_{j \geq K}|B_{ij}|^2\Bigr)^{1/2}.
\]
By Parseval, $\sum_{j \geq K} |S_{i,j+K}^+(\delta,L)|^2 \leq \|T_\delta e_i\|^2 \leq 1$. Then:
\[
  \|B\|_{\mathrm{HS}}^2 \leq K\cdot\Bigl(\sum_{p,m}\frac{2\log p}{p^{m/2}}\Bigr)^2 \leq K\cdot\bigl(4\sqrt{\lambda}\bigr)^2 = 16K\lambda.
\]
So $\|B\|_{\mathrm{op}} \leq 4\sqrt{K\lambda} = 4\sqrt{K}\,e^{L/2}$.

\textbf{This is $O(e^{L/2})$, not $O(L)$.} The numerical scaling $\|B\|/L \in [0.2, 1.6]$ holds only on the tested range $\lambda \in [30, 500]$; asymptotically, $\|B\|$ grows as $e^{L/2}$.
\end{proof}

\begin{remark}[Coupling asymmetry---an honest assessment]
\label{rem:coupling-revisited}
The Hilbert--Schmidt bound gives $\|B\|_{\mathrm{op}} = O(\sqrt{\lambda})$---the \emph{same order} as the dominant diagonal $W_{11}^+$. The coupling correction is therefore \emph{not negligible}.

One might hope that the coupling ``cancels'' in the gap $\Delta = \lambda_1^+ - \lambda_1^-$, since both sectors share the same high-mode structure. However, numerical tests reveal a \textbf{coupling asymmetry}: $\|B^-\|_{\mathrm{op}}$ is systematically $30$--$40\%$ larger than $\|B^+\|_{\mathrm{op}}$ (e.g., $\|B^-\| = 4.47$ vs.\ $\|B^+\| = 3.40$ at $\lambda = 100$). This means the odd eigenvalue is pushed \emph{further down} by coupling than the even eigenvalue, which \emph{reduces} the gap. The high block is parity-neutral in $\lambda_1(C)$ (difference $< 0.02$) but \emph{not} in $\|B\|$.

Consequently, the gap cannot be rigorously reduced to the diagonal difference $W_{11}^+ - W_{00}^-$ alone. The coupling asymmetry is a genuine analytical obstruction that our block-bound argument does not resolve. Specifically:
\begin{itemize}
  \item The Rayleigh quotient gives $\lambda_1^+ \leq W_{11}^+$ (upper bound on the even eigenvalue).
  \item For the gap, one needs a \emph{lower} bound on $\lambda_1^-$ (the odd eigenvalue).
  \item The Schur complement gives $\lambda_1^- \geq W_{00}^- - \|B^-\|^2/(\lambda_1(C^-) - W_{00}^-)$, but since $\|B^-\| = O(\sqrt{\lambda})$ and the denominator is also $O(\sqrt{\lambda})$, this bound is $W_{00}^- - O(\sqrt{\lambda})$---not useful, as the correction is the same order as the leading term.
\end{itemize}
\end{remark}

\subsection{The Certifier Meta-Theorem}
\label{sec:meta-theorem}

The key to closing the proof is that the \emph{certified gap} grows monotonically with $\lambda$, not because of abstract operator convergence, but because the Galerkin eigenvalues of the finite blocks scale cleanly.

\begin{definition}[Certified gap]
For fixed $k$ (even block size) and $N_0$ (odd block size), define
\begin{align*}
  \ell_{\cos}^{(k)}(\lambda) &\coloneqq \lambda_{\min}\!\bigl(\mathrm{QW}_{\cos}(\lambda)\big|_{\text{modes }0\ldots k-1}\bigr), \\
  \ell_{\sin}^{(N_0)}(\lambda) &\coloneqq \lambda_{\min}\!\bigl(\mathrm{QW}_{\sin}(\lambda)\big|_{\text{modes }0\ldots N_0-1}\bigr), \\
  \mathrm{cert\_gap}(\lambda) &\coloneqq \ell_{\cos}^{(k)}(\lambda) - \ell_{\sin}^{(N_0)}(\lambda).
\end{align*}
By Cauchy interlacing, $\ell_{\cos}^{(k)}(\lambda) \geq \lambda_1^+(\lambda)$ and $\ell_{\sin}^{(N_0)}(\lambda) \geq \lambda_1^-(\lambda)$, so
\[
  \mathrm{cert\_gap}(\lambda) < 0 \;\implies\; \lambda_1^+(\lambda) < \lambda_1^-(\lambda) \quad\text{(even dominance).}
\]
\end{definition}

\begin{proposition}[Asymptotic scaling of finite-block eigenvalues]
\label{prop:scaling}
For $k = 4$ and any fixed $N_0$, with primes summed up to $\lambda$:
\begin{align}
  \ell_{\cos}^{(4)}(\lambda) / \sqrt{\lambda} &\;\to\; -c_{\cos} \quad (\lambda \to \infty), \label{eq:ccos-limit} \\
  \ell_{\sin}^{(N_0)}(\lambda) / \sqrt{\lambda} &\;\to\; -c_{\sin}^{(N_0)} \quad (\lambda \to \infty), \label{eq:csin-limit}
\end{align}
where $c_{\cos}$ and $c_{\sin}^{(N_0)}$ are positive constants determined by limit matrices.
\end{proposition}

\noindent\textbf{Numerical evidence.} With primes correctly summed up to~$\lambda$:

\begin{center}
\renewcommand{\arraystretch}{1.15}
\begin{tabular}{r|r|r|r|r}
$\lambda$ & $\#$primes & $\ell_{\cos}^{(4)}/\sqrt{\lambda}$ & $\ell_{\sin}^{(4)}/\sqrt{\lambda}$ & gap$/\sqrt{\lambda}$ \\\hline
100  & 25  & $-1.102$ & $-0.807$ & $-0.295$ \\
200  & 46  & $-1.343$ & $-0.973$ & $-0.370$ \\
500  & 95  & $-1.527$ & $-1.115$ & $-0.412$ \\
1000 & 168 & $-1.649$ & $-1.211$ & $-0.438$ \\
2000 & 303 & $-1.761$ & $-1.303$ & $-0.458$
\end{tabular}
\end{center}

\noindent All three ratios converge monotonically. The gap ratio stabilizes at $\approx -0.46$, confirming $c_{\cos} > c_{\sin}$.

\begin{theorem}[Meta-theorem: certifier termination]
\label{thm:meta}
Assume:
\begin{enumerate}[label=\textup{(M\arabic*)}]
  \item\label{M1} The limits~\eqref{eq:ccos-limit}--\eqref{eq:csin-limit} exist with $c_{\cos} > c_{\sin}^{(N_0)}$.
  \item\label{M2} The tail correction $\mathrm{tail}(\lambda; N_0) \coloneqq \ell_{\sin}^{(N_0)}(\lambda) - \lambda_1^-(\lambda)$ satisfies $\mathrm{tail}(\lambda; N_0) = O(\log\lambda)$.
  \item\label{M3} The interval-arithmetic evaluator produces certified bounds with width $E(\lambda) = O(\lambda^{1/2-\varepsilon})$ for some $\varepsilon > 0$.
\end{enumerate}
Then there exists $\lambda_0$ such that for all $\lambda \geq \lambda_0$:
\begin{enumerate}[label=\textup{(\roman*)}]
  \item $\mathrm{cert\_gap}(\lambda) \leq -(c_{\cos} - c_{\sin}^{(N_0)})\sqrt{\lambda} + O(\log\lambda) < 0$,
  \item the interval certifier terminates and outputs ``even dominance verified,''
  \item $\lambda_1^+(\lambda) < \lambda_1^-(\lambda)$ (even dominance holds).
\end{enumerate}
In particular, there exists a sequence $\lambda_n \to \infty$ along which the hypotheses of Connes' Theorem~6.1 are rigorously verified, and by Fact~6.4 and Hurwitz's theorem, all zeros of $\Xi$ lie on the real line.
\end{theorem}

\begin{proof}
By \ref{M1}: $\mathrm{cert\_gap}(\lambda) = -(c_{\cos} - c_{\sin}^{(N_0)} + o(1))\sqrt{\lambda}$. Since $c_{\cos} - c_{\sin}^{(N_0)} > 0$, the gap is eventually negative and $|\mathrm{cert\_gap}| \to \infty$.

By \ref{M3}: the interval evaluator produces an interval $I(\lambda)$ containing $\mathrm{cert\_gap}(\lambda)$ with $\mathrm{rad}(I) \leq E(\lambda) = O(\lambda^{1/2-\varepsilon})$. Since $|\mathrm{cert\_gap}(\lambda)| \sim (c_{\cos}-c_{\sin})\sqrt{\lambda}$, we have $|\mathrm{cert\_gap}| / E(\lambda) \to \infty$. Hence $\sup I(\lambda) < 0$ for all $\lambda \geq \lambda_0$, and the certifier decides ``gap $< 0$'' in finitely many refinement steps.

By \ref{M2}: $\lambda_1^-(\lambda) \geq \ell_{\sin}^{(N_0)}(\lambda) - O(\log\lambda) > \ell_{\cos}^{(k)}(\lambda) \geq \lambda_1^+(\lambda)$, establishing even dominance.

Connes' Theorem~6.1 then gives: for each such $\lambda$, the entire function associated to the ground-state eigenfunction has only real zeros. Fact~6.4 gives $k_\lambda \to \Xi$ locally uniformly. Hurwitz's theorem transfers the real-zero property to the limit $\Xi$. Hence all non-trivial zeros of $\zeta$ lie on the critical line.
\end{proof}

\subsection{Status of the assumptions}

\begin{center}
\renewcommand{\arraystretch}{1.3}
\begin{tabular}{l|l|l}
\textbf{Assumption} & \textbf{Status} & \textbf{What is needed} \\\hline
\ref{M1}: $c_{\cos} > c_{\sin}^{(N_0)}$ & Near-proved & PNT limit of matrix entries + Shift Parity \\
& & on the $m=1$ prime sum (\Cref{cor:diagonal-parity}) \\
\ref{M2}: $\mathrm{tail} = O(\log\lambda)$ & Proved (for fixed $N_0$) & Cauchy interlacing + Mertens \\
\ref{M3}: $E(\lambda) = O(\lambda^{1/2-\varepsilon})$ & Standard & Interval arithmetic with $O(N_0^2)$ entries, \\
& & each computed in $O(\pi(\lambda))$ operations
\end{tabular}
\end{center}

\noindent The critical assumption is \ref{M1}. It asserts that the $4 \times 4$ even-block eigenvalue is \emph{asymptotically more negative} than the $N_0 \times N_0$ odd-block eigenvalue, after normalizing by $\sqrt{\lambda}$. The mechanism is the Shift Parity identity (\Cref{cor:diagonal-parity}): the cosine overlap $S^+(\delta, L) < S^-(\delta, L)$ for $\delta \in (0, L)$, which makes the even matrix entries more negative. The challenge is to prove that this pointwise inequality on the overlaps propagates through the eigenvalue computation (which involves the full matrix, not just diagonal entries).

Numerical evidence strongly supports \ref{M1}: the ratio $\mathrm{gap}/\sqrt{\lambda}$ is monotonically increasing in magnitude from $-0.295$ ($\lambda = 100$) to $-0.458$ ($\lambda = 2000$), with no sign of reversal.

\begin{remark}[Concrete path to \ref{M1}]
Define the \emph{limit matrices} $C_{\cos} \coloneqq \lim_{\lambda \to \infty} \mathrm{QW}_{\cos}^{(k)}(\lambda) / \sqrt{\lambda}$ and $C_{\sin} \coloneqq \lim_{\lambda \to \infty} \mathrm{QW}_{\sin}^{(N_0)}(\lambda) / \sqrt{\lambda}$. The entries of these matrices are Stieltjes integrals:
\[
  (C_{\pm})_{nm} = \lim_{\lambda \to \infty} \frac{2}{\sqrt{\lambda}} \sum_{p \leq \lambda} \frac{\log p}{\sqrt{p}}\,S_{nm}^{\pm}(\log p,\, \log\lambda).
\]
By PNT and partial summation with the smooth weight $S_{nm}^{\pm}(\log p, L)$, these limits exist and equal explicit integrals over $[0,1]$ (after the substitution $u = \log p / L$). The Shift Parity identity $S_{nm}^+ - S_{nm}^- = -\sin(\cdot)/\pi$ then gives $(C_{\cos})_{nm} < (C_{\sin})_{nm}$ entry-wise for the leading terms, from which $\lambda_{\min}(C_{\cos}) < \lambda_{\min}(C_{\sin})$ follows by Weyl's inequality---provided the off-diagonal corrections are controlled.

A complete proof of \ref{M1} along these lines requires:
\begin{enumerate}
  \item Existence and explicit computation of $C_{\cos}$ and $C_{\sin}$ (PNT + dominated convergence).
  \item $\lambda_{\min}(C_{\cos}) < \lambda_{\min}(C_{\sin})$ (Weyl + Shift Parity on the matrix level).
\end{enumerate}
Both are accessible with standard analytic number theory tools.
\end{remark}


\bibliographystyle{plain}
\bibliography{fst-rh-references}

\end{document}
